\clearpage
\section{Why one revision and two revisions print the same interval}

On 25 of the 26 rows the two interval columns are identical, in the printed band and in the
population. The only row where they differ is \texttt{mediator}. The reason is structural, and it is the same
fact the figure is built on seen from another angle.

The interval over $r$ revisions is a union over knowledge states, so it grows with $r$ and it is
bounded above by the interval over the whole equivalence class. The class interval is a ceiling that
no amount of retracting can pass, because retracting all of the knowledge is exactly what the class
already assumes. So the two columns can only differ on a row where the first revision has not yet
reached that ceiling.

\begin{center}\small
\begin{tabular}{@{}rp{12.6cm}@{}}
\toprule
rows & what the second revision adds, and why \\
\midrule
9 & The class interval is a point. The effect is identified by the data alone, no asserted claim
bears on it, and every budget returns the same sampling interval. Seven of these are declared
queries whose claims all sit outside the treatment's component \\
\addlinespace
16 & The first revision already returns the whole class, to four decimals. The ceiling is reached at
$r = 1$, so $r = 2$ has nothing left to add \\
\addlinespace
1 & mediator. The first revision covers 57\,\% of the class and the second saturates it.
This is the only row in the table where the second column carries information \\
\bottomrule
\end{tabular}
\end{center}

The mechanism behind the sixteen is the breakdown radius itself. The extreme of the class is
attained by some knowledge state, and reaching it requires retracting the smallest set of claims
that reaches it. That set is a single claim on 141 of the 157 queries with an exact certificate, and
on every one of those the union at $r = 1$ already contains the extreme. A second retraction can
only add states that were already inside.

\subsection{The same count on all 303 rows}

The exhaustive table lets the question be asked of the whole population rather than of the
twenty-six.

\begin{center}\small
\begin{tabular}{@{}rp{12.6cm}@{}}
\toprule
rows & of the 303 queries that commit to an adjustment set \\
\midrule
281 & print the same interval at one and at two revisions, in the band and in the population,
which is 92.7\,\% \\
76 & have a class interval that is a point: no asserted claim bears on the query at any budget \\
205 & reach the whole class at the first revision \\
22 & have a first revision strictly inside the class, and on 20 of those the second revision
widens the interval \\
\bottomrule
\end{tabular}
\end{center}

So the second column carries information on 20 queries in 303. Where those 20 sit is the part worth
noticing: they are on insurance, Polzer 2012, Kampen 2014, sachs, Schipf 2010 and mediator. Every
one of those networks except mediator is excluded from the printed T3 by the real-coefficient rule.
The column looks empty in T3 because T3 drops the networks on which it is full.

\begin{verdict}
\textbf{What to do with the column.} Replace \emph{two revisions} with \emph{no knowledge at all},
the class interval. On the sixteen saturated rows the reader then sees two identical intervals side
by side, and that identity is the finding rather than a redundancy. On mediator they differ,
and the difference is legible for the first time. The exhaustive table below prints all three
columns so the substitution can be checked.
\end{verdict}

\clearpage
\section{The exhaustive table}

Both selection rules removed. Every query under the recovering set whose analyst graph commits to an
adjustment set: 303 rows over all 27 swept networks, with the four fitted-coefficient networks and
the 23 semi-synthetic ones in separate blocks and a provenance column, so the rule that was dropped
is visible rather than applied. A bullet marks the 26 rows the printed T3 selects.

Every cell that the printed T3 also shows was checked against it before this table was written, and
all 26 rows agree cell by cell.

Two queries appear twice, once as a declared pair and once as an independently sampled one:
Thoemmes 2013, x on y, and mediator, X on Y. They are the same query reached by two routes, and both
rows are kept.

\begin{landscape}
\begingroup
\fontsize{6.3}{7.3}\selectfont
\setlength{\tabcolsep}{1.5pt}
\renewcommand{\arraystretch}{1.04}
\begin{longtable}{@{}llp{2.15cm}ccrl cccc c l@{}}
\caption*{\textbf{T3-X, the exhaustive table.} Every query under the recovering set whose analyst graph commits to an adjustment set: 303 rows over 27 networks. Neither rule of the printed T3 is applied here, so networks with parameters we drew appear beside the four with coefficients fitted to real data (\emph{param.}), and no per-network cap is used. A bullet after the network name marks the 26 rows the printed T3 selects. \emph{loc.} is the number of asserted claims inside the treatment's own chain component. The \emph{no knowledge} column is the interval over the whole equivalence class, which is the ceiling every other interval is bounded by. Estimates at $n=20\,000$, first seed; the two revision columns are bands at that seed, the class column is the population.}\\
\toprule
network & param. & query & $|K|$ & loc. & true & estimate [95\,\%] & one revision & two revisions & no knowledge & $r_0$ & $r_{\mathrm{val}}$ & first breaking claim \\
\midrule
\endfirsthead
\toprule
network & param. & query & $|K|$ & loc. & true & estimate [95\,\%] & one revision & two revisions & no knowledge & $r_0$ & $r_{\mathrm{val}}$ & first breaking claim \\
\midrule
\endhead
\midrule \multicolumn{13}{r@{}}{\itshape continues on the next page}\\
\endfoot
\bottomrule
\endlastfoot
\multicolumn{13}{@{}l}{\itshape Queries the source papers declare.}\\[1pt]
Acid\_1996$\bullet$ & semi-s. & x3 $\to$ x15 & 1 & 0 & $-2.169$ & $-2.176$ [$-2.203$, $-2.149$] & [$-2.226$, $-2.147$] & [$-2.226$, $-2.147$] & [$-2.169$, $-2.169$] & never & never & -- \\
Didelez\_2010$\bullet$ & semi-s. & HRT $\to$ TCI & 1 & 0 & $-0.586$ & $-0.586$ [$-0.599$, $-0.572$] & [$-0.604$, $-0.568$] & [$-0.604$, $-0.568$] & [$-0.586$, $-0.586$] & never & never & -- \\
Kampen\_2014$\bullet$ & semi-s. & SUS $\to$ EGC & 3 & 0 & $-0.643$ & $-0.647$ [$-0.663$, $-0.630$] & [$-0.663$, $-0.630$] & [$-0.663$, $-0.630$] & [$-0.643$, $-0.643$] & never & never & -- \\
Polzer\_2012$\bullet$ & semi-s. & Tooth\allowbreak{}Loss $\to$ Mortality & 5 & 0 & $\phantom{-}1.193$ & $\phantom{-}1.190$ [$\phantom{-}1.183$, $\phantom{-}1.197$] & [$\phantom{-}1.128$, $\phantom{-}1.208$] & [$\phantom{-}1.128$, $\phantom{-}1.208$] & [$\phantom{-}1.193$, $\phantom{-}1.193$] & never & $>$3 & -- \\
Sebastiani\_2005$\bullet$ & semi-s. & EDN1.3 $\to$ EDNI1.7 & 7 & 0 & $-1.212$ & $-1.209$ [$-1.219$, $-1.199$] & [$-1.234$, $-1.179$] & [$-1.234$, $-1.179$] & [$-1.212$, $-1.212$] & never & $>$3 & -- \\
Shrier\_2008$\bullet$ & semi-s. & Warm\allowbreak{}Up\allowbreak{}Exercises $\to$ Injury & 2 & 0 & $-1.351$ & $-1.348$ [$-1.358$, $-1.339$] & [$-1.431$, $-1.228$] & [$-1.431$, $-1.228$] & [$-1.351$, $-1.351$] & never & never & -- \\
Thoemmes\_2013$\bullet$ & semi-s. & x $\to$ y & 1 & 1 & $-0.923$ & $-0.913$ [$-0.928$, $-0.898$] & [$-1.002$, $-0.871$] & [$-1.002$, $-0.871$] & [$-0.923$, $-0.923$] & never & never & -- \\
mediator$\bullet$ & semi-s. & X $\to$ Y & 3 & 3 & $-2.603$ & $-2.600$ [$-2.623$, $-2.576$] & [$-2.623$, $-1.109$] & [$-2.623$, $\phantom{-}0.000$] & [$-2.603$, $\phantom{-}0.000$] & 2 & 1 & Z$\to$I \,/\, Z$\to$X \\
paths$\bullet$ & semi-s. & E $\to$ D & 1 & 0 & $\phantom{-}1.194$ & $\phantom{-}1.193$ [$\phantom{-}1.188$, $\phantom{-}1.197$] & [$\phantom{-}1.169$, $\phantom{-}1.216$] & [$\phantom{-}1.169$, $\phantom{-}1.216$] & [$\phantom{-}1.194$, $\phantom{-}1.194$] & never & never & -- \\
\addlinespace
\multicolumn{13}{@{}l}{\itshape Sampled queries, networks with coefficients fitted to real data.}\\[1pt]
arth150$\bullet$ & fitted & 100 $\to$ 245 & 29 & 3 & $\phantom{-}1.535$ & $\phantom{-}1.534$ [$\phantom{-}1.527$, $\phantom{-}1.542$] & [$\phantom{-}0.000$, $\phantom{-}1.545$] & [$\phantom{-}0.000$, $\phantom{-}1.545$] & [$\phantom{-}0.000$, $\phantom{-}1.535$] & 1 & 1 & 100$\to$245 \\
arth150$\bullet$ & fitted & 81 $\to$ 464 & 29 & 10 & $\phantom{-}0.256$ & $\phantom{-}0.184$ [$\phantom{-}0.118$, $\phantom{-}0.250$] & [$\phantom{-}0.097$, $\phantom{-}0.265$] & [$\phantom{-}0.097$, $\phantom{-}0.265$] & [$\phantom{-}0.256$, $\phantom{-}0.256$] & $>$3 & $\geq$2 & -- \\
ecoli70$\bullet$ & fitted & asn\allowbreak{}A $\to$ ace\allowbreak{}B & 10 & 1 & $\phantom{-}0.547$ & $\phantom{-}0.542$ [$\phantom{-}0.511$, $\phantom{-}0.572$] & [$-0.011$, $\phantom{-}0.572$] & [$-0.011$, $\phantom{-}0.572$] & [$\phantom{-}0.000$, $\phantom{-}0.547$] & 1 & 1 & asnA$\to$icdA \\
ecoli70$\bullet$ & fitted & b1191 $\to$ yfa\allowbreak{}D & 10 & 1 & $-0.098$ & $-0.097$ [$-0.105$, $-0.089$] & [$-0.114$, $\phantom{-}0.029$] & [$-0.114$, $\phantom{-}0.029$] & [$-0.098$, $-0.000$] & 1 & 1 & b1191$\to$fixC \\
ecoli70$\bullet$ & fitted & csp\allowbreak{}G $\to$ hup\allowbreak{}B & 10 & 7 & $\phantom{-}0.257$ & $\phantom{-}0.265$ [$\phantom{-}0.246$, $\phantom{-}0.283$] & [$\phantom{-}0.003$, $\phantom{-}0.308$] & [$\phantom{-}0.003$, $\phantom{-}0.308$] & [$\phantom{-}0.000$, $\phantom{-}0.257$] & 1 & 1 & -- \\
ecoli70$\bullet$ & fitted & eut\allowbreak{}G $\to$ lac\allowbreak{}Y & 10 & 1 & $\phantom{-}0.309$ & $\phantom{-}0.309$ [$\phantom{-}0.279$, $\phantom{-}0.338$] & [$\phantom{-}0.277$, $\phantom{-}0.377$] & [$\phantom{-}0.277$, $\phantom{-}0.377$] & [$\phantom{-}0.309$, $\phantom{-}0.354$] & never & 1 & -- \\
ecoli70 & fitted & icd\allowbreak{}A $\to$ ace\allowbreak{}B & 10 & 1 & $\phantom{-}1.046$ & $\phantom{-}1.049$ [$\phantom{-}1.046$, $\phantom{-}1.052$] & [$\phantom{-}1.043$, $\phantom{-}1.058$] & [$\phantom{-}1.043$, $\phantom{-}1.058$] & [$\phantom{-}1.046$, $\phantom{-}1.046$] & never & $>$3 & -- \\
ecoli70$\bullet$ & fitted & psp\allowbreak{}B $\to$ nmp\allowbreak{}C & 10 & 7 & $-0.108$ & $-0.100$ [$-0.118$, $-0.082$] & [$-0.531$, $\phantom{-}0.026$] & [$-0.531$, $\phantom{-}0.026$] & [$-0.525$, $\phantom{-}0.000$] & 1 & 1 & pspB$\to$pspA \\
ecoli70 & fitted & suc\allowbreak{}A $\to$ tna\allowbreak{}A & 10 & 1 & $\phantom{-}0.111$ & $\phantom{-}0.113$ [$\phantom{-}0.109$, $\phantom{-}0.116$] & [$\phantom{-}0.086$, $\phantom{-}0.127$] & [$\phantom{-}0.086$, $\phantom{-}0.127$] & [$\phantom{-}0.111$, $\phantom{-}0.111$] & never & $>$3 & -- \\
ecoli70 & fitted & yed\allowbreak{}E $\to$ b1963 & 10 & 7 & $\phantom{-}0.351$ & $\phantom{-}0.352$ [$\phantom{-}0.340$, $\phantom{-}0.363$] & [$\phantom{-}0.259$, $\phantom{-}0.421$] & [$\phantom{-}0.259$, $\phantom{-}0.421$] & [$\phantom{-}0.351$, $\phantom{-}0.351$] & $>$3 & $>$3 & -- \\
magic-irri$\bullet$ & fitted & G1378 $\to$ SUB & 10 & 1 & $-0.035$ & $-0.031$ [$-0.044$, $-0.017$] & [$-0.048$, $\phantom{-}0.021$] & [$-0.048$, $\phantom{-}0.021$] & [$-0.035$, $\phantom{-}0.000$] & 1 & 1 & G1378$\to$G1311 \\
magic-irri & fitted & G2670 $\to$ YLD & 10 & 1 & $-0.082$ & $-0.086$ [$-0.095$, $-0.076$] & [$-0.099$, $-0.077$] & [$-0.099$, $-0.077$] & [$-0.082$, $-0.082$] & never & $>$3 & -- \\
magic-irri & fitted & G3209 $\to$ G3106 & 10 & 1 & $-0.189$ & $-0.203$ [$-0.219$, $-0.186$] & [$-0.218$, $-0.171$] & [$-0.218$, $-0.171$] & [$-0.189$, $-0.189$] & never & $>$3 & -- \\
magic-irri$\bullet$ & fitted & G3212 $\to$ FT & 10 & 2 & $\phantom{-}0.004$ & $-0.025$ [$-0.070$, $\phantom{-}0.020$] & [$-0.143$, $\phantom{-}0.008$] & [$-0.143$, $\phantom{-}0.008$] & [$-0.042$, $\phantom{-}0.004$] & 1 & 1 & -- \\
magic-irri & fitted & G3219 $\to$ G3102 & 10 & 1 & $-0.072$ & $-0.088$ [$-0.101$, $-0.075$] & [$-0.102$, $-0.075$] & [$-0.102$, $-0.075$] & [$-0.072$, $-0.072$] & never & $>$3 & -- \\
magic-irri & fitted & G3219 $\to$ HT & 10 & 1 & $\phantom{-}0.070$ & $\phantom{-}0.030$ [$-0.082$, $\phantom{-}0.142$] & [$-0.137$, $\phantom{-}0.176$] & [$-0.137$, $\phantom{-}0.176$] & [$\phantom{-}0.004$, $\phantom{-}0.070$] & never & 1 & -- \\
magic-irri & fitted & G3823 $\to$ HT & 10 & 1 & $\phantom{-}0.359$ & $\phantom{-}0.318$ [$\phantom{-}0.208$, $\phantom{-}0.428$] & [$\phantom{-}0.234$, $\phantom{-}0.485$] & [$\phantom{-}0.234$, $\phantom{-}0.485$] & [$\phantom{-}0.359$, $\phantom{-}0.359$] & never & $>$3 & -- \\
magic-irri$\bullet$ & fitted & G3925 $\to$ GTEMP & 10 & 1 & $\phantom{-}0.000$ & $\phantom{-}0.001$ [$-0.034$, $\phantom{-}0.036$] & [$-0.131$, $\phantom{-}0.038$] & [$-0.131$, $\phantom{-}0.038$] & [$\phantom{-}0.000$, $\phantom{-}0.000$] & 1 & 1 & -- \\
magic-irri$\bullet$ & fitted & G4145 $\to$ G3222 & 10 & 2 & $-0.020$ & $-0.015$ [$-0.025$, $-0.004$] & [$-0.026$, $\phantom{-}0.026$] & [$-0.026$, $\phantom{-}0.026$] & [$-0.020$, $-0.000$] & 1 & 1 & -- \\
magic-irri$\bullet$ & fitted & G4156 $\to$ CHALK & 10 & 2 & $\phantom{-}0.404$ & $\phantom{-}0.501$ [$\phantom{-}0.361$, $\phantom{-}0.641$] & [$-0.154$, $\phantom{-}0.664$] & [$-0.154$, $\phantom{-}0.664$] & [$\phantom{-}0.000$, $\phantom{-}0.404$] & 1 & 1 & G4156$\to$G4145 \\
magic-irri & fitted & G4156 $\to$ G3862 & 10 & 2 & $\phantom{-}0.000$ & $\phantom{-}0.001$ [$-0.007$, $\phantom{-}0.010$] & [$-0.007$, $\phantom{-}0.017$] & [$-0.007$, $\phantom{-}0.017$] & [$\phantom{-}0.000$, $\phantom{-}0.000$] & 1 & 1 & -- \\
magic-irri & fitted & G4382 $\to$ G3098 & 10 & 1 & $\phantom{-}0.008$ & $\phantom{-}0.006$ [$-0.008$, $\phantom{-}0.020$] & [$-0.018$, $\phantom{-}0.020$] & [$-0.018$, $\phantom{-}0.020$] & [$\phantom{-}0.000$, $\phantom{-}0.008$] & 1 & 1 & -- \\
magic-irri & fitted & G4382 $\to$ GW & 10 & 1 & $\phantom{-}0.002$ & $\phantom{-}0.002$ [$\phantom{-}0.000$, $\phantom{-}0.005$] & [$-0.003$, $\phantom{-}0.005$] & [$-0.003$, $\phantom{-}0.005$] & [$\phantom{-}0.000$, $\phantom{-}0.002$] & 1 & 1 & -- \\
magic-niab$\bullet$ & fitted & G1217 $\to$ FT & 8 & 1 & $\phantom{-}0.068$ & $\phantom{-}0.015$ [$-0.029$, $\phantom{-}0.060$] & [$-0.031$, $\phantom{-}0.131$] & [$-0.031$, $\phantom{-}0.131$] & [$\phantom{-}0.068$, $\phantom{-}0.122$] & never & 1 & -- \\
magic-niab$\bullet$ & fitted & G1217 $\to$ YLD & 8 & 1 & $\phantom{-}0.009$ & $\phantom{-}0.011$ [$\phantom{-}0.003$, $\phantom{-}0.018$] & [$\phantom{-}0.004$, $\phantom{-}0.021$] & [$\phantom{-}0.004$, $\phantom{-}0.021$] & [$\phantom{-}0.007$, $\phantom{-}0.009$] & never & 1 & -- \\
magic-niab & fitted & G1276 $\to$ G1789 & 8 & 1 & $\phantom{-}0.039$ & $\phantom{-}0.032$ [$\phantom{-}0.022$, $\phantom{-}0.042$] & [$\phantom{-}0.020$, $\phantom{-}0.042$] & [$\phantom{-}0.020$, $\phantom{-}0.042$] & [$\phantom{-}0.039$, $\phantom{-}0.039$] & never & $>$3 & -- \\
magic-niab & fitted & G1276 $\to$ YR.FIELD & 8 & 1 & $\phantom{-}0.015$ & $\phantom{-}0.018$ [$\phantom{-}0.012$, $\phantom{-}0.024$] & [$\phantom{-}0.012$, $\phantom{-}0.041$] & [$\phantom{-}0.012$, $\phantom{-}0.041$] & [$\phantom{-}0.015$, $\phantom{-}0.030$] & never & 1 & -- \\
magic-niab & fitted & G1750 $\to$ YLD & 8 & 1 & $-0.011$ & $-0.006$ [$-0.015$, $\phantom{-}0.003$] & [$-0.014$, $\phantom{-}0.008$] & [$-0.014$, $\phantom{-}0.008$] & [$-0.011$, $-0.011$] & never & $>$3 & -- \\
magic-niab$\bullet$ & fitted & G1853 $\to$ G43 & 8 & 1 & $-0.039$ & $-0.028$ [$-0.038$, $-0.017$] & [$-0.038$, $\phantom{-}0.020$] & [$-0.038$, $\phantom{-}0.020$] & [$-0.039$, $\phantom{-}0.000$] & 1 & 1 & G1853$\to$G311 \\
magic-niab & fitted & G1896 $\to$ FUS & 8 & 1 & $-0.242$ & $-0.243$ [$-0.255$, $-0.231$] & [$-0.258$, $-0.226$] & [$-0.258$, $-0.226$] & [$-0.242$, $-0.234$] & never & 1 & -- \\
magic-niab & fitted & G2208 $\to$ MIL & 8 & 1 & $\phantom{-}0.424$ & $\phantom{-}0.412$ [$\phantom{-}0.375$, $\phantom{-}0.449$] & [$\phantom{-}0.372$, $\phantom{-}0.480$] & [$\phantom{-}0.372$, $\phantom{-}0.480$] & [$\phantom{-}0.424$, $\phantom{-}0.454$] & never & 1 & -- \\
magic-niab & fitted & G257 $\to$ YR.FIELD & 8 & 1 & $-0.082$ & $-0.080$ [$-0.087$, $-0.073$] & [$-0.096$, $-0.072$] & [$-0.096$, $-0.072$] & [$-0.090$, $-0.081$] & never & 1 & -- \\
magic-niab & fitted & G2835 $\to$ G1800 & 8 & 2 & $\phantom{-}0.124$ & $\phantom{-}0.122$ [$\phantom{-}0.110$, $\phantom{-}0.133$] & [$\phantom{-}0.110$, $\phantom{-}0.135$] & [$\phantom{-}0.110$, $\phantom{-}0.135$] & [$\phantom{-}0.124$, $\phantom{-}0.124$] & never & $>$3 & -- \\
magic-niab & fitted & G2953 $\to$ HT & 8 & 1 & $-1.026$ & $-1.014$ [$-1.082$, $-0.945$] & [$-1.393$, $-0.935$] & [$-1.393$, $-0.935$] & [$-1.349$, $-1.026$] & never & 1 & -- \\
magic-niab$\bullet$ & fitted & G311 $\to$ G43 & 8 & 1 & $-0.255$ & $-0.251$ [$-0.259$, $-0.244$] & [$-0.260$, $\phantom{-}0.000$] & [$-0.260$, $\phantom{-}0.000$] & [$-0.255$, $\phantom{-}0.000$] & 1 & 1 & G1853$\to$G311 \\
magic-niab$\bullet$ & fitted & G418 $\to$ FUS & 8 & 2 & $\phantom{-}0.003$ & $-0.003$ [$-0.015$, $\phantom{-}0.010$] & [$-0.022$, $\phantom{-}0.011$] & [$-0.022$, $\phantom{-}0.011$] & [$\phantom{-}0.000$, $\phantom{-}0.003$] & 1 & 1 & -- \\
magic-niab & fitted & G418 $\to$ YLD & 8 & 2 & $\phantom{-}0.043$ & $\phantom{-}0.045$ [$\phantom{-}0.037$, $\phantom{-}0.052$] & [$\phantom{-}0.028$, $\phantom{-}0.054$] & [$\phantom{-}0.028$, $\phantom{-}0.054$] & [$\phantom{-}0.036$, $\phantom{-}0.045$] & never & 1 & -- \\
\addlinespace
\multicolumn{13}{@{}l}{\itshape Sampled queries, networks whose parameters we drew (structure from a published source).}\\[1pt]
Acid\_1996 & semi-s. & x1 $\to$ x10 & 1 & 1 & $-2.725$ & $-2.686$ [$-2.735$, $-2.637$] & [$-2.749$, $-1.315$] & [$-2.749$, $-1.315$] & [$-2.725$, $-1.433$] & never & 1 & -- \\
Acid\_1996 & semi-s. & x1 $\to$ x11 & 1 & 1 & $\phantom{-}3.064$ & $\phantom{-}3.026$ [$\phantom{-}2.988$, $\phantom{-}3.064$] & [$\phantom{-}0.772$, $\phantom{-}3.072$] & [$\phantom{-}0.772$, $\phantom{-}3.072$] & [$\phantom{-}0.847$, $\phantom{-}3.064$] & never & 1 & -- \\
Acid\_1996 & semi-s. & x1 $\to$ x12 & 1 & 1 & $\phantom{-}6.934$ & $\phantom{-}6.844$ [$\phantom{-}6.744$, $\phantom{-}6.944$] & [$\phantom{-}2.684$, $\phantom{-}6.972$] & [$\phantom{-}2.684$, $\phantom{-}6.972$] & [$\phantom{-}2.918$, $\phantom{-}6.934$] & never & 1 & -- \\
Acid\_1996 & semi-s. & x1 $\to$ x13 & 1 & 1 & $-9.327$ & $-9.211$ [$-9.347$, $-9.076$] & [$-9.384$, $-3.611$] & [$-9.384$, $-3.611$] & [$-9.327$, $-3.925$] & never & 1 & -- \\
Acid\_1996 & semi-s. & x1 $\to$ x14 & 1 & 1 & $\phantom{-}0.674$ & $\phantom{-}0.661$ [$\phantom{-}0.636$, $\phantom{-}0.685$] & [$\phantom{-}0.569$, $\phantom{-}0.684$] & [$\phantom{-}0.569$, $\phantom{-}0.684$] & [$\phantom{-}0.674$, $\phantom{-}0.674$] & never & never & -- \\
Acid\_1996 & semi-s. & x1 $\to$ x15 & 1 & 1 & $\phantom{-}2.446$ & $\phantom{-}2.412$ [$\phantom{-}2.365$, $\phantom{-}2.458$] & [$\phantom{-}1.150$, $\phantom{-}2.458$] & [$\phantom{-}1.150$, $\phantom{-}2.458$] & [$\phantom{-}1.287$, $\phantom{-}2.446$] & never & 1 & -- \\
Acid\_1996 & semi-s. & x1 $\to$ x16 & 1 & 1 & $-0.565$ & $-0.545$ [$-0.569$, $-0.520$] & [$-0.568$, $-0.462$] & [$-0.568$, $-0.462$] & [$-0.565$, $-0.565$] & never & never & -- \\
Acid\_1996 & semi-s. & x1 $\to$ x17 & 1 & 1 & $-3.277$ & $-3.213$ [$-3.283$, $-3.143$] & [$-3.283$, $-1.778$] & [$-3.283$, $-1.778$] & [$-3.277$, $-2.032$] & never & 1 & -- \\
Acid\_1996 & semi-s. & x1 $\to$ x18 & 1 & 1 & $\phantom{-}2.450$ & $\phantom{-}2.409$ [$\phantom{-}2.361$, $\phantom{-}2.458$] & [$\phantom{-}1.140$, $\phantom{-}2.458$] & [$\phantom{-}1.140$, $\phantom{-}2.458$] & [$\phantom{-}1.288$, $\phantom{-}2.450$] & never & 1 & -- \\
Acid\_1996 & semi-s. & x1 $\to$ x3 & 1 & 1 & $-0.593$ & $-0.585$ [$-0.599$, $-0.571$] & [$-0.606$, $-0.547$] & [$-0.606$, $-0.547$] & [$-0.593$, $-0.593$] & never & never & -- \\
Acid\_1996 & semi-s. & x1 $\to$ x4 & 1 & 1 & $-1.382$ & $-1.366$ [$-1.380$, $-1.352$] & [$-1.380$, $\phantom{-}0.000$] & [$-1.380$, $\phantom{-}0.000$] & [$-1.382$, $\phantom{-}0.000$] & 1 & 1 & x1$\to$x4 \\
Acid\_1996 & semi-s. & x1 $\to$ x5 & 1 & 1 & $-2.127$ & $-2.105$ [$-2.130$, $-2.080$] & [$-2.135$, $-0.537$] & [$-2.135$, $-0.537$] & [$-2.127$, $-0.588$] & never & 1 & -- \\
Acid\_1996 & semi-s. & x1 $\to$ x6 & 1 & 1 & $\phantom{-}0.675$ & $\phantom{-}0.659$ [$\phantom{-}0.639$, $\phantom{-}0.680$] & [$\phantom{-}0.597$, $\phantom{-}0.687$] & [$\phantom{-}0.597$, $\phantom{-}0.687$] & [$\phantom{-}0.675$, $\phantom{-}0.675$] & never & never & -- \\
Acid\_1996 & semi-s. & x1 $\to$ x7 & 1 & 1 & $\phantom{-}2.536$ & $\phantom{-}2.499$ [$\phantom{-}2.457$, $\phantom{-}2.541$] & [$\phantom{-}1.206$, $\phantom{-}2.554$] & [$\phantom{-}1.206$, $\phantom{-}2.554$] & [$\phantom{-}1.334$, $\phantom{-}2.536$] & never & 1 & -- \\
Acid\_1996 & semi-s. & x1 $\to$ x9 & 1 & 1 & $\phantom{-}2.702$ & $\phantom{-}2.665$ [$\phantom{-}2.618$, $\phantom{-}2.711$] & [$\phantom{-}1.296$, $\phantom{-}2.725$] & [$\phantom{-}1.296$, $\phantom{-}2.725$] & [$\phantom{-}1.421$, $\phantom{-}2.702$] & never & 1 & -- \\
Acid\_1996 & semi-s. & x4 $\to$ x10 & 1 & 1 & $\phantom{-}0.935$ & $\phantom{-}0.922$ [$\phantom{-}0.905$, $\phantom{-}0.939$] & [$\phantom{-}0.895$, $\phantom{-}1.648$] & [$\phantom{-}0.895$, $\phantom{-}1.648$] & [$\phantom{-}0.935$, $\phantom{-}1.616$] & never & 1 & -- \\
Acid\_1996 & semi-s. & x4 $\to$ x11 & 1 & 1 & $-1.605$ & $-1.601$ [$-1.616$, $-1.586$] & [$-2.029$, $-1.579$] & [$-2.029$, $-1.579$] & [$-2.007$, $-1.605$] & never & 1 & -- \\
Acid\_1996 & semi-s. & x4 $\to$ x12 & 1 & 1 & $-2.907$ & $-2.887$ [$-2.922$, $-2.853$] & [$-4.360$, $-2.832$] & [$-4.360$, $-2.832$] & [$-4.293$, $-2.907$] & never & 1 & -- \\
Acid\_1996 & semi-s. & x4 $\to$ x13 & 1 & 1 & $\phantom{-}3.910$ & $\phantom{-}3.889$ [$\phantom{-}3.842$, $\phantom{-}3.936$] & [$\phantom{-}3.812$, $\phantom{-}5.869$] & [$\phantom{-}3.812$, $\phantom{-}5.869$] & [$\phantom{-}3.910$, $\phantom{-}5.774$] & never & 1 & -- \\
Didelez\_2010 & semi-s. & Age $\to$ HRT & 1 & 1 & $\phantom{-}0.577$ & $\phantom{-}0.570$ [$\phantom{-}0.543$, $\phantom{-}0.596$] & [$-0.772$, $\phantom{-}0.596$] & [$-0.772$, $\phantom{-}0.596$] & [$-0.729$, $\phantom{-}0.577$] & 1 & 1 & Age$\to$Occ \\
Didelez\_2010 & semi-s. & Age $\to$ Occ & 1 & 1 & $\phantom{-}1.219$ & $\phantom{-}1.212$ [$\phantom{-}1.198$, $\phantom{-}1.226$] & [$\phantom{-}0.000$, $\phantom{-}1.226$] & [$\phantom{-}0.000$, $\phantom{-}1.226$] & [$\phantom{-}0.000$, $\phantom{-}1.219$] & 1 & 1 & Age$\to$Occ \\
Didelez\_2010 & semi-s. & Age $\to$ S & 1 & 1 & $\phantom{-}1.976$ & $\phantom{-}1.977$ [$\phantom{-}1.951$, $\phantom{-}2.003$] & [$\phantom{-}1.112$, $\phantom{-}2.003$] & [$\phantom{-}1.112$, $\phantom{-}2.003$] & [$\phantom{-}1.162$, $\phantom{-}1.976$] & never & 1 & -- \\
Didelez\_2010 & semi-s. & Age $\to$ Smo & 1 & 1 & $\phantom{-}1.054$ & $\phantom{-}1.046$ [$\phantom{-}1.028$, $\phantom{-}1.065$] & [$-0.025$, $\phantom{-}1.065$] & [$-0.025$, $\phantom{-}1.065$] & [$-0.000$, $\phantom{-}1.054$] & 1 & 1 & Age$\to$Occ \\
Didelez\_2010 & semi-s. & Age $\to$ TCI & 1 & 1 & $-1.991$ & $-1.984$ [$-2.018$, $-1.950$] & [$-2.018$, $-0.689$] & [$-2.018$, $-0.689$] & [$-1.991$, $-0.753$] & never & 1 & -- \\
Didelez\_2010 & semi-s. & Age $\to$ Thist & 1 & 1 & $\phantom{-}0.960$ & $\phantom{-}0.959$ [$\phantom{-}0.941$, $\phantom{-}0.976$] & [$-0.008$, $\phantom{-}0.976$] & [$-0.008$, $\phantom{-}0.976$] & [$-0.000$, $\phantom{-}0.960$] & 1 & 1 & Age$\to$Occ \\
Didelez\_2010 & semi-s. & Occ $\to$ HRT & 1 & 1 & $\phantom{-}1.071$ & $\phantom{-}1.078$ [$\phantom{-}1.056$, $\phantom{-}1.100$] & [$-0.363$, $\phantom{-}1.100$] & [$-0.363$, $\phantom{-}1.100$] & [$-0.357$, $\phantom{-}1.071$] & 1 & 1 & Age$\to$Occ \\
Didelez\_2010 & semi-s. & Occ $\to$ S & 1 & 1 & $\phantom{-}0.668$ & $\phantom{-}0.682$ [$\phantom{-}0.657$, $\phantom{-}0.706$] & [$\phantom{-}0.217$, $\phantom{-}1.601$] & [$\phantom{-}0.217$, $\phantom{-}1.601$] & [$\phantom{-}0.218$, $\phantom{-}1.589$] & never & 1 & -- \\
Didelez\_2010 & semi-s. & Occ $\to$ Smo & 1 & 1 & $\phantom{-}0.865$ & $\phantom{-}0.864$ [$\phantom{-}0.855$, $\phantom{-}0.873$] & [$\phantom{-}0.000$, $\phantom{-}0.879$] & [$\phantom{-}0.000$, $\phantom{-}0.879$] & [$\phantom{-}0.000$, $\phantom{-}0.865$] & 1 & 1 & Age$\to$Occ \\
Didelez\_2010 & semi-s. & Occ $\to$ TCI & 1 & 1 & $-1.016$ & $-1.028$ [$-1.059$, $-0.997$] & [$-1.938$, $\phantom{-}0.167$] & [$-1.938$, $\phantom{-}0.167$] & [$-1.920$, $\phantom{-}0.166$] & 1 & 1 & Age$\to$Occ \\
Didelez\_2010 & semi-s. & Occ $\to$ Thist & 1 & 1 & $\phantom{-}0.788$ & $\phantom{-}0.786$ [$\phantom{-}0.777$, $\phantom{-}0.795$] & [$\phantom{-}0.000$, $\phantom{-}0.795$] & [$\phantom{-}0.000$, $\phantom{-}0.795$] & [$\phantom{-}0.000$, $\phantom{-}0.788$] & 1 & 1 & Age$\to$Occ \\
Didelez\_2010 & semi-s. & Smo $\to$ HRT & 1 & 1 & $\phantom{-}1.238$ & $\phantom{-}1.236$ [$\phantom{-}1.226$, $\phantom{-}1.247$] & [$\phantom{-}0.960$, $\phantom{-}1.245$] & [$\phantom{-}0.960$, $\phantom{-}1.245$] & [$\phantom{-}0.970$, $\phantom{-}1.238$] & never & 1 & -- \\
Didelez\_2010 & semi-s. & Smo $\to$ S & 1 & 1 & $\phantom{-}1.179$ & $\phantom{-}1.175$ [$\phantom{-}1.161$, $\phantom{-}1.189$] & [$\phantom{-}1.145$, $\phantom{-}1.358$] & [$\phantom{-}1.145$, $\phantom{-}1.358$] & [$\phantom{-}1.179$, $\phantom{-}1.342$] & never & 1 & -- \\
Didelez\_2010 & semi-s. & Smo $\to$ TCI & 1 & 1 & $-1.793$ & $-1.792$ [$-1.805$, $-1.779$] & [$-1.798$, $-1.657$] & [$-1.798$, $-1.657$] & [$-1.793$, $-1.668$] & never & 1 & -- \\
Kampen\_2014 & semi-s. & AFF $\to$ CDR & 3 & 3 & $\phantom{-}1.443$ & $\phantom{-}1.442$ [$\phantom{-}1.434$, $\phantom{-}1.450$] & [$\phantom{-}0.000$, $\phantom{-}1.735$] & [$\phantom{-}0.000$, $\phantom{-}1.735$] & [$\phantom{-}0.000$, $\phantom{-}1.724$] & 1 & 1 & AIS$\to$AFF \\
Kampen\_2014 & semi-s. & AFF $\to$ FTW & 3 & 3 & $\phantom{-}1.351$ & $\phantom{-}1.387$ [$\phantom{-}1.346$, $\phantom{-}1.427$] & [$\phantom{-}0.382$, $\phantom{-}1.919$] & [$\phantom{-}0.382$, $\phantom{-}1.919$] & [$\phantom{-}0.412$, $\phantom{-}1.844$] & never & 1 & -- \\
Kampen\_2014 & semi-s. & AIS $\to$ ALN & 3 & 3 & $-0.738$ & $-0.744$ [$-0.759$, $-0.728$] & [$-0.759$, $\phantom{-}0.035$] & [$-0.759$, $\phantom{-}0.575$] & [$-0.738$, $\phantom{-}0.546$] & 1 & 1 & AIS$\to$AFF \\
Kampen\_2014 & semi-s. & AIS $\to$ DET & 3 & 3 & $-4.948$ & $-4.940$ [$-4.997$, $-4.882$] & [$-4.997$, $-1.031$] & [$-4.997$, $\phantom{-}0.952$] & [$-4.948$, $\phantom{-}0.883$] & 2 & 1 & AIS$\to$AFF \,/\, SAN$\to$AIS \\
Kampen\_2014 & semi-s. & AIS $\to$ HOS & 3 & 3 & $-4.166$ & $-4.164$ [$-4.211$, $-4.117$] & [$-4.211$, $-2.994$] & [$-4.211$, $-2.220$] & [$-4.166$, $-2.283$] & never & 1 & -- \\
Kampen\_2014 & semi-s. & AIS $\to$ SUS & 3 & 3 & $\phantom{-}1.887$ & $\phantom{-}1.888$ [$\phantom{-}1.864$, $\phantom{-}1.913$] & [$\phantom{-}0.896$, $\phantom{-}1.913$] & [$\phantom{-}0.224$, $\phantom{-}1.913$] & [$\phantom{-}0.255$, $\phantom{-}1.887$] & never & 1 & -- \\
Kampen\_2014 & semi-s. & ALN $\to$ APA & 3 & 3 & $\phantom{-}0.583$ & $\phantom{-}0.592$ [$\phantom{-}0.578$, $\phantom{-}0.605$] & [$\phantom{-}0.000$, $\phantom{-}1.305$] & [$\phantom{-}0.000$, $\phantom{-}1.305$] & [$\phantom{-}0.000$, $\phantom{-}1.264$] & 1 & 1 & ALN$\to$APA \\
Kampen\_2014 & semi-s. & ALN $\to$ EGC & 3 & 3 & $-0.024$ & $-0.024$ [$-0.038$, $-0.011$] & [$-0.568$, $\phantom{-}0.521$] & [$-0.568$, $\phantom{-}0.521$] & [$-0.547$, $\phantom{-}0.506$] & 1 & 1 & AIS$\to$AFF \\
Kampen\_2014 & semi-s. & ALN $\to$ PER & 3 & 3 & $\phantom{-}1.473$ & $\phantom{-}1.478$ [$\phantom{-}1.468$, $\phantom{-}1.488$] & [$\phantom{-}0.000$, $\phantom{-}1.497$] & [$\phantom{-}0.000$, $\phantom{-}1.497$] & [$\phantom{-}0.000$, $\phantom{-}1.473$] & 1 & 1 & AIS$\to$AFF \\
Kampen\_2014 & semi-s. & CDR $\to$ DET & 3 & 3 & $\phantom{-}1.168$ & $\phantom{-}1.173$ [$\phantom{-}1.166$, $\phantom{-}1.180$] & [$\phantom{-}1.146$, $\phantom{-}2.184$] & [$\phantom{-}1.146$, $\phantom{-}2.184$] & [$\phantom{-}1.168$, $\phantom{-}2.155$] & never & 1 & -- \\
Polzer\_2012 & semi-s. & Alcohol $\to$ Lipids & 5 & 3 & $-2.856$ & $-2.839$ [$-2.881$, $-2.797$] & [$-3.374$, $-2.706$] & [$-3.374$, $-2.545$] & [$-3.348$, $-2.625$] & never & 1 & -- \\
Polzer\_2012 & semi-s. & Alcohol $\to$ Periodontitis & 5 & 3 & $-0.060$ & $-0.042$ [$-0.091$, $\phantom{-}0.007$] & [$-2.030$, $\phantom{-}0.766$] & [$-2.111$, $\phantom{-}0.766$] & [$-2.061$, $\phantom{-}0.706$] & 1 & 1 & -- \\
Polzer\_2012 & semi-s. & Alcohol $\to$ Sport & 5 & 3 & $\phantom{-}0.797$ & $\phantom{-}0.800$ [$\phantom{-}0.783$, $\phantom{-}0.817$] & [$-0.018$, $\phantom{-}0.817$] & [$-0.018$, $\phantom{-}0.817$] & [$\phantom{-}0.000$, $\phantom{-}0.799$] & 1 & 1 & Alcohol$\to$Smoking \\
Polzer\_2012 & semi-s. & Caries $\to$ Mortality & 5 & 3 & $\phantom{-}1.094$ & $\phantom{-}1.092$ [$\phantom{-}1.071$, $\phantom{-}1.113$] & [$-0.877$, $\phantom{-}1.162$] & [$-0.877$, $\phantom{-}1.162$] & [$-0.806$, $\phantom{-}1.094$] & 1 & 1 & Psychosocial$\to$Caries \\
Polzer\_2012 & semi-s. & Caries $\to$ Tooth\allowbreak{}Loss & 5 & 3 & $\phantom{-}0.917$ & $\phantom{-}0.916$ [$\phantom{-}0.902$, $\phantom{-}0.930$] & [$-0.515$, $\phantom{-}1.002$] & [$-0.515$, $\phantom{-}1.002$] & [$-0.476$, $\phantom{-}0.917$] & 1 & 1 & Psychosocial$\to$Caries \\
Polzer\_2012 & semi-s. & Diabetes $\to$ Mortality & 5 & 2 & $\phantom{-}0.919$ & $\phantom{-}0.922$ [$\phantom{-}0.886$, $\phantom{-}0.959$] & [$-0.121$, $\phantom{-}2.293$] & [$-0.121$, $\phantom{-}2.318$] & [$-0.097$, $\phantom{-}2.241$] & 1 & 1 & Obesity$\to$Diabetes \\
Polzer\_2012 & semi-s. & Diabetes $\to$ Tooth\allowbreak{}Loss & 5 & 2 & $\phantom{-}1.973$ & $\phantom{-}1.979$ [$\phantom{-}1.954$, $\phantom{-}2.003$] & [$\phantom{-}1.211$, $\phantom{-}2.023$] & [$\phantom{-}1.211$, $\phantom{-}2.053$] & [$\phantom{-}1.225$, $\phantom{-}1.989$] & never & 1 & -- \\
Polzer\_2012 & semi-s. & Hypertension $\to$ Mortality & 5 & 2 & $\phantom{-}0.985$ & $\phantom{-}0.990$ [$\phantom{-}0.976$, $\phantom{-}1.004$] & [$-0.111$, $\phantom{-}1.427$] & [$-0.111$, $\phantom{-}1.427$] & [$-0.085$, $\phantom{-}1.395$] & 1 & 1 & Diabetes$\to$Hypertension \\
Polzer\_2012 & semi-s. & Obesity $\to$ Lipids & 5 & 2 & $\phantom{-}0.109$ & $\phantom{-}0.117$ [$\phantom{-}0.095$, $\phantom{-}0.139$] & [$\phantom{-}0.095$, $\phantom{-}1.296$] & [$\phantom{-}0.095$, $\phantom{-}1.296$] & [$\phantom{-}0.109$, $\phantom{-}1.273$] & never & 1 & -- \\
Polzer\_2012 & semi-s. & Obesity $\to$ Tooth\allowbreak{}Loss & 5 & 2 & $\phantom{-}0.285$ & $\phantom{-}0.294$ [$\phantom{-}0.258$, $\phantom{-}0.331$] & [$-1.553$, $\phantom{-}0.338$] & [$-1.553$, $\phantom{-}0.338$] & [$-1.521$, $\phantom{-}0.285$] & 1 & 1 & Diabetes$\to$Hypertension \\
Polzer\_2012 & semi-s. & Psychosocial $\to$ Diabetes & 5 & 3 & $\phantom{-}0.601$ & $\phantom{-}0.615$ [$\phantom{-}0.579$, $\phantom{-}0.650$] & [$\phantom{-}0.553$, $\phantom{-}2.300$] & [$\phantom{-}0.553$, $\phantom{-}2.300$] & [$\phantom{-}0.601$, $\phantom{-}2.258$] & never & 1 & -- \\
Polzer\_2012 & semi-s. & Psychosocial $\to$ Mortality & 5 & 3 & $-2.726$ & $-2.743$ [$-2.835$, $-2.651$] & [$-3.910$, $\phantom{-}0.687$] & [$-3.910$, $\phantom{-}0.687$] & [$-3.800$, $\phantom{-}0.592$] & 1 & 1 & Alcohol$\to$Smoking \\
Polzer\_2012 & semi-s. & Psychosocial $\to$ Smoking & 5 & 3 & $\phantom{-}1.312$ & $\phantom{-}1.315$ [$\phantom{-}1.294$, $\phantom{-}1.336$] & [$\phantom{-}0.000$, $\phantom{-}1.354$] & [$\phantom{-}0.000$, $\phantom{-}1.354$] & [$\phantom{-}0.000$, $\phantom{-}1.312$] & 1 & 1 & Alcohol$\to$Smoking \\
Polzer\_2012 & semi-s. & Smoking $\to$ Hypertension & 5 & 3 & $\phantom{-}3.949$ & $\phantom{-}3.939$ [$\phantom{-}3.896$, $\phantom{-}3.982$] & [$\phantom{-}0.132$, $\phantom{-}3.982$] & [$\phantom{-}0.132$, $\phantom{-}3.982$] & [$\phantom{-}0.174$, $\phantom{-}3.949$] & never & 1 & -- \\
Polzer\_2012 & semi-s. & Smoking $\to$ Obesity & 5 & 3 & $-1.541$ & $-1.542$ [$-1.561$, $-1.523$] & [$-1.561$, $\phantom{-}0.135$] & [$-1.561$, $\phantom{-}0.135$] & [$-1.541$, $\phantom{-}0.123$] & 1 & 1 & Alcohol$\to$Smoking \\
Schipf\_2010 & semi-s. & A $\to$ PA & 3 & 3 & $-1.264$ & $-1.259$ [$-1.273$, $-1.245$] & [$-1.278$, $\phantom{-}0.000$] & [$-1.278$, $\phantom{-}0.000$] & [$-1.264$, $\phantom{-}0.000$] & 1 & 1 & A$\to$PA \\
Schipf\_2010 & semi-s. & A $\to$ S & 3 & 3 & $\phantom{-}1.235$ & $\phantom{-}1.236$ [$\phantom{-}1.222$, $\phantom{-}1.250$] & [$\phantom{-}0.000$, $\phantom{-}1.255$] & [$\phantom{-}0.000$, $\phantom{-}1.255$] & [$\phantom{-}0.000$, $\phantom{-}1.235$] & 1 & 1 & A$\to$S \\
Schipf\_2010 & semi-s. & A $\to$ T2DM & 3 & 3 & $\phantom{-}0.363$ & $\phantom{-}0.368$ [$\phantom{-}0.338$, $\phantom{-}0.399$] & [$-1.924$, $\phantom{-}1.553$] & [$-1.924$, $\phantom{-}1.553$] & [$-1.931$, $\phantom{-}1.514$] & 1 & 1 & A$\to$PA \\
Schipf\_2010 & semi-s. & A $\to$ TT & 3 & 3 & $-2.942$ & $-2.949$ [$-2.978$, $-2.921$] & [$-3.907$, $-1.660$] & [$-3.907$, $-1.660$] & [$-3.878$, $-1.685$] & never & 1 & -- \\
Schipf\_2010 & semi-s. & A $\to$ U & 3 & 3 & $-1.564$ & $-1.574$ [$-1.597$, $-1.550$] & [$-1.604$, $\phantom{-}0.870$] & [$-1.604$, $\phantom{-}0.870$] & [$-1.564$, $\phantom{-}0.867$] & 1 & 1 & A$\to$PA \\
Schipf\_2010 & semi-s. & A $\to$ WC & 3 & 3 & $-2.428$ & $-2.433$ [$-2.452$, $-2.413$] & [$-2.452$, $\phantom{-}0.000$] & [$-2.452$, $\phantom{-}0.000$] & [$-2.428$, $\phantom{-}0.000$] & 1 & 1 & A$\to$PA \\
Schipf\_2010 & semi-s. & PA $\to$ T2DM & 3 & 3 & $-1.451$ & $-1.433$ [$-1.452$, $-1.414$] & [$-1.454$, $-0.720$] & [$-1.454$, $-0.720$] & [$-1.451$, $-0.735$] & never & 1 & -- \\
Schipf\_2010 & semi-s. & PA $\to$ TT & 3 & 3 & $-0.160$ & $-0.142$ [$-0.166$, $-0.117$] & [$-0.167$, $\phantom{-}1.398$] & [$-0.167$, $\phantom{-}1.398$] & [$-0.160$, $\phantom{-}1.370$] & 1 & 1 & A$\to$PA \\
Schipf\_2010 & semi-s. & PA $\to$ U & 3 & 3 & $\phantom{-}0.979$ & $\phantom{-}0.977$ [$\phantom{-}0.957$, $\phantom{-}0.996$] & [$-0.028$, $\phantom{-}1.155$] & [$-0.043$, $\phantom{-}1.155$] & [$-0.021$, $\phantom{-}1.138$] & 1 & 1 & PA$\to$WC \\
Schipf\_2010 & semi-s. & PA $\to$ WC & 3 & 3 & $\phantom{-}0.977$ & $\phantom{-}0.988$ [$\phantom{-}0.974$, $\phantom{-}1.002$] & [$\phantom{-}0.000$, $\phantom{-}1.574$] & [$\phantom{-}0.000$, $\phantom{-}1.574$] & [$\phantom{-}0.000$, $\phantom{-}1.557$] & 1 & 1 & PA$\to$WC \\
Schipf\_2010 & semi-s. & S $\to$ T2DM & 3 & 3 & $-0.932$ & $-0.924$ [$-0.936$, $-0.912$] & [$-0.950$, $-0.173$] & [$-0.950$, $-0.173$] & [$-0.932$, $-0.192$] & never & 1 & -- \\
Schipf\_2010 & semi-s. & S $\to$ TT & 3 & 3 & $-1.018$ & $-1.008$ [$-1.022$, $-0.994$] & [$-1.852$, $-0.987$] & [$-1.852$, $-0.987$] & [$-1.842$, $-1.018$] & never & 1 & -- \\
Sebastiani\_2005 & semi-s. & ANXA2.8 $\to$ ANXA2.7 & 7 & 3 & $\phantom{-}1.155$ & $\phantom{-}1.151$ [$\phantom{-}1.130$, $\phantom{-}1.173$] & [$\phantom{-}1.104$, $\phantom{-}1.188$] & [$\phantom{-}1.104$, $\phantom{-}1.188$] & [$\phantom{-}1.155$, $\phantom{-}1.155$] & never & $>$3 & -- \\
Sebastiani\_2005 & semi-s. & ANXA2.8 $\to$ CSF2.4 & 7 & 3 & $\phantom{-}0.782$ & $\phantom{-}0.796$ [$\phantom{-}0.774$, $\phantom{-}0.818$] & [$-0.070$, $\phantom{-}0.862$] & [$-0.070$, $\phantom{-}0.862$] & [$\phantom{-}0.000$, $\phantom{-}0.782$] & 1 & 1 & ANXA2.8$\to$CAT \\
Sebastiani\_2005 & semi-s. & ANXA2.8 $\to$ SELP.12 & 7 & 3 & $\phantom{-}1.197$ & $\phantom{-}1.193$ [$\phantom{-}1.164$, $\phantom{-}1.221$] & [$-0.057$, $\phantom{-}1.262$] & [$-0.057$, $\phantom{-}1.262$] & [$\phantom{-}0.000$, $\phantom{-}1.197$] & 1 & 1 & ANXA2.8$\to$CAT \\
Sebastiani\_2005 & semi-s. & BMP6 $\to$ BMP6.13 & 7 & 3 & $\phantom{-}0.515$ & $\phantom{-}0.516$ [$\phantom{-}0.506$, $\phantom{-}0.527$] & [$\phantom{-}0.501$, $\phantom{-}0.602$] & [$\phantom{-}0.501$, $\phantom{-}0.602$] & [$\phantom{-}0.515$, $\phantom{-}0.515$] & never & $>$3 & -- \\
Sebastiani\_2005 & semi-s. & CAT $\to$ CSF2.4 & 7 & 3 & $\phantom{-}0.530$ & $\phantom{-}0.541$ [$\phantom{-}0.529$, $\phantom{-}0.553$] & [$\phantom{-}0.519$, $\phantom{-}0.575$] & [$\phantom{-}0.519$, $\phantom{-}0.575$] & [$\phantom{-}0.530$, $\phantom{-}0.530$] & never & $>$3 & -- \\
Sebastiani\_2005 & semi-s. & CAT $\to$ SELP.12 & 7 & 3 & $\phantom{-}0.812$ & $\phantom{-}0.806$ [$\phantom{-}0.791$, $\phantom{-}0.820$] & [$\phantom{-}0.782$, $\phantom{-}0.836$] & [$\phantom{-}0.782$, $\phantom{-}0.836$] & [$\phantom{-}0.812$, $\phantom{-}0.812$] & never & $>$3 & -- \\
Sebastiani\_2005 & semi-s. & EDN1.9 $\to$ ANXA2.12 & 7 & 1 & $-0.489$ & $-0.483$ [$-0.502$, $-0.464$] & [$-0.518$, $-0.452$] & [$-0.518$, $-0.452$] & [$-0.489$, $-0.489$] & never & $>$3 & -- \\
Sebastiani\_2005 & semi-s. & EDNI1.6 $\to$ EDN1.10 & 7 & 3 & $\phantom{-}0.668$ & $\phantom{-}0.674$ [$\phantom{-}0.666$, $\phantom{-}0.683$] & [$\phantom{-}0.666$, $\phantom{-}0.706$] & [$\phantom{-}0.666$, $\phantom{-}0.706$] & [$\phantom{-}0.668$, $\phantom{-}0.668$] & never & $>$3 & -- \\
Sebastiani\_2005 & semi-s. & Stroke $\to$ BMP6.9 & 7 & 1 & $\phantom{-}1.756$ & $\phantom{-}1.785$ [$\phantom{-}1.757$, $\phantom{-}1.813$] & [$\phantom{-}1.734$, $\phantom{-}1.830$] & [$\phantom{-}1.734$, $\phantom{-}1.830$] & [$\phantom{-}1.756$, $\phantom{-}1.756$] & never & $>$3 & -- \\
Sebastiani\_2005 & semi-s. & Stroke $\to$ SELP.17 & 7 & 1 & $\phantom{-}0.185$ & $\phantom{-}0.190$ [$\phantom{-}0.169$, $\phantom{-}0.210$] & [$\phantom{-}0.117$, $\phantom{-}0.211$] & [$\phantom{-}0.117$, $\phantom{-}0.211$] & [$\phantom{-}0.185$, $\phantom{-}0.185$] & never & $>$3 & -- \\
Shrier\_2008 & semi-s. & Coach $\to$ Fitness\allowbreak{}Level & 2 & 1 & $-1.012$ & $-1.018$ [$-1.032$, $-1.005$] & [$-1.043$, $-0.994$] & [$-1.043$, $-0.994$] & [$-1.012$, $-1.012$] & never & never & -- \\
Shrier\_2008 & semi-s. & Coach $\to$ Intra\allowbreak{}Game\allowbreak{}Proprioception & 2 & 1 & $-2.335$ & $-2.334$ [$-2.368$, $-2.301$] & [$-2.374$, $-1.433$] & [$-2.374$, $-1.433$] & [$-2.335$, $-1.491$] & never & 1 & -- \\
Shrier\_2008 & semi-s. & Coach $\to$ Pre\allowbreak{}Game\allowbreak{}Proprioception & 2 & 1 & $-1.235$ & $-1.241$ [$-1.263$, $-1.219$] & [$-1.286$, $-1.213$] & [$-1.286$, $-1.213$] & [$-1.235$, $-1.235$] & never & never & -- \\
Shrier\_2008 & semi-s. & Coach $\to$ Team\allowbreak{}Motivation & 2 & 1 & $-1.069$ & $-1.071$ [$-1.085$, $-1.057$] & [$-1.085$, $\phantom{-}0.000$] & [$-1.085$, $\phantom{-}0.000$] & [$-1.069$, $\phantom{-}0.000$] & 1 & 1 & Coach$\to$TeamMotivation \\
Shrier\_2008 & semi-s. & Connective\allowbreak{}Tissue\allowbreak{}Disorder $\to$ Injury & 2 & 1 & $-3.044$ & $-3.022$ [$-3.057$, $-2.986$] & [$-5.041$, $-2.884$] & [$-5.041$, $-2.884$] & [$-5.042$, $-3.044$] & never & 1 & -- \\
Shrier\_2008 & semi-s. & Connective\allowbreak{}Tissue\allowbreak{}Disorder $\to$ Neuromuscular\allowbreak{}Fatigue & 2 & 1 & $\phantom{-}1.297$ & $\phantom{-}1.294$ [$\phantom{-}1.280$, $\phantom{-}1.308$] & [$\phantom{-}1.263$, $\phantom{-}2.188$] & [$\phantom{-}1.263$, $\phantom{-}2.188$] & [$\phantom{-}1.297$, $\phantom{-}2.182$] & never & 1 & -- \\
Shrier\_2008 & semi-s. & Connective\allowbreak{}Tissue\allowbreak{}Disorder $\to$ Tissue\allowbreak{}Weakness & 2 & 1 & $\phantom{-}0.838$ & $\phantom{-}0.830$ [$\phantom{-}0.822$, $\phantom{-}0.838$] & [$\phantom{-}0.000$, $\phantom{-}0.843$] & [$\phantom{-}0.000$, $\phantom{-}0.843$] & [$\phantom{-}0.000$, $\phantom{-}0.838$] & 1 & 1 & Genetics$\to$ConnectiveTissueDisorder \\
Shrier\_2008 & semi-s. & Genetics $\to$ Connective\allowbreak{}Tissue\allowbreak{}Disorder & 2 & 1 & $\phantom{-}1.365$ & $\phantom{-}1.368$ [$\phantom{-}1.354$, $\phantom{-}1.382$] & [$\phantom{-}0.000$, $\phantom{-}1.382$] & [$\phantom{-}0.000$, $\phantom{-}1.382$] & [$\phantom{-}0.000$, $\phantom{-}1.365$] & 1 & 1 & Genetics$\to$ConnectiveTissueDisorder \\
Shrier\_2008 & semi-s. & Genetics $\to$ Injury & 2 & 1 & $-8.345$ & $-8.344$ [$-8.416$, $-8.272$] & [$-8.425$, $-4.137$] & [$-8.425$, $-4.137$] & [$-8.345$, $-4.191$] & never & 1 & -- \\
Shrier\_2008 & semi-s. & Genetics $\to$ Neuromuscular\allowbreak{}Fatigue & 2 & 1 & $\phantom{-}3.626$ & $\phantom{-}3.632$ [$\phantom{-}3.607$, $\phantom{-}3.657$] & [$\phantom{-}1.840$, $\phantom{-}3.658$] & [$\phantom{-}1.840$, $\phantom{-}3.658$] & [$\phantom{-}1.855$, $\phantom{-}3.626$] & never & 1 & -- \\
Shrier\_2008 & semi-s. & Genetics $\to$ Tissue\allowbreak{}Weakness & 2 & 1 & $\phantom{-}1.144$ & $\phantom{-}1.137$ [$\phantom{-}1.119$, $\phantom{-}1.155$] & [$-0.021$, $\phantom{-}1.155$] & [$-0.021$, $\phantom{-}1.155$] & [$-0.000$, $\phantom{-}1.144$] & 1 & 1 & Genetics$\to$ConnectiveTissueDisorder \\
Shrier\_2008 & semi-s. & Team\allowbreak{}Motivation $\to$ Injury & 2 & 1 & $\phantom{-}1.056$ & $\phantom{-}1.042$ [$\phantom{-}1.019$, $\phantom{-}1.064$] & [$\phantom{-}0.855$, $\phantom{-}2.478$] & [$\phantom{-}0.855$, $\phantom{-}2.478$] & [$\phantom{-}1.056$, $\phantom{-}2.437$] & never & 1 & -- \\
Shrier\_2008 & semi-s. & Team\allowbreak{}Motivation $\to$ Previous\allowbreak{}Injury & 2 & 1 & $\phantom{-}0.648$ & $\phantom{-}0.647$ [$\phantom{-}0.638$, $\phantom{-}0.657$] & [$\phantom{-}0.633$, $\phantom{-}0.678$] & [$\phantom{-}0.633$, $\phantom{-}0.678$] & [$\phantom{-}0.648$, $\phantom{-}0.648$] & never & never & -- \\
Thoemmes\_2013 & semi-s. & e0 $\to$ s1 & 1 & 1 & $\phantom{-}0.813$ & $\phantom{-}0.802$ [$\phantom{-}0.786$, $\phantom{-}0.818$] & [$-0.058$, $\phantom{-}0.812$] & [$-0.058$, $\phantom{-}0.812$] & [$\phantom{-}0.000$, $\phantom{-}0.813$] & 1 & 1 & e0$\to$x \\
Thoemmes\_2013 & semi-s. & e0 $\to$ s2 & 1 & 1 & $\phantom{-}0.509$ & $\phantom{-}0.496$ [$\phantom{-}0.479$, $\phantom{-}0.513$] & [$-0.097$, $\phantom{-}0.503$] & [$-0.097$, $\phantom{-}0.503$] & [$\phantom{-}0.000$, $\phantom{-}0.509$] & 1 & 1 & e0$\to$x \\
Thoemmes\_2013 & semi-s. & e0 $\to$ s3 & 1 & 1 & $\phantom{-}0.652$ & $\phantom{-}0.637$ [$\phantom{-}0.611$, $\phantom{-}0.663$] & [$-0.120$, $\phantom{-}0.647$] & [$-0.120$, $\phantom{-}0.647$] & [$-0.000$, $\phantom{-}0.652$] & 1 & 1 & e0$\to$x \\
Thoemmes\_2013 & semi-s. & e0 $\to$ x & 1 & 1 & $-1.368$ & $-1.356$ [$-1.370$, $-1.342$] & [$-1.370$, $\phantom{-}0.000$] & [$-1.370$, $\phantom{-}0.000$] & [$-1.368$, $\phantom{-}0.000$] & 1 & 1 & e0$\to$x \\
Thoemmes\_2013 & semi-s. & e0 $\to$ y & 1 & 1 & $\phantom{-}1.262$ & $\phantom{-}1.239$ [$\phantom{-}1.210$, $\phantom{-}1.267$] & [$-0.188$, $\phantom{-}1.243$] & [$-0.188$, $\phantom{-}1.243$] & [$-0.000$, $\phantom{-}1.262$] & 1 & 1 & e0$\to$x \\
Thoemmes\_2013 & semi-s. & e0 $\to$ z & 1 & 1 & $-0.577$ & $-0.578$ [$-0.596$, $-0.560$] & [$-0.590$, $\phantom{-}0.045$] & [$-0.590$, $\phantom{-}0.045$] & [$-0.577$, $\phantom{-}0.000$] & 1 & 1 & e0$\to$x \\
Thoemmes\_2013 & semi-s. & e0 $\to$ z2 & 1 & 1 & $-0.275$ & $-0.278$ [$-0.295$, $-0.262$] & [$-0.287$, $\phantom{-}0.037$] & [$-0.287$, $\phantom{-}0.037$] & [$-0.275$, $\phantom{-}0.000$] & 1 & 1 & e0$\to$x \\
Thoemmes\_2013 & semi-s. & e0 $\to$ z3 & 1 & 1 & $\phantom{-}1.380$ & $\phantom{-}1.355$ [$\phantom{-}1.321$, $\phantom{-}1.389$] & [$-0.189$, $\phantom{-}1.361$] & [$-0.189$, $\phantom{-}1.361$] & [$-0.000$, $\phantom{-}1.380$] & 1 & 1 & e0$\to$x \\
Thoemmes\_2013 & semi-s. & x $\to$ s1 & 1 & 1 & $-0.594$ & $-0.591$ [$-0.600$, $-0.583$] & [$-0.620$, $-0.576$] & [$-0.620$, $-0.576$] & [$-0.594$, $-0.594$] & never & never & -- \\
Thoemmes\_2013 & semi-s. & x $\to$ s2 & 1 & 1 & $-0.373$ & $-0.365$ [$-0.375$, $-0.355$] & [$-0.416$, $-0.348$] & [$-0.416$, $-0.348$] & [$-0.373$, $-0.373$] & never & never & -- \\
Thoemmes\_2013 & semi-s. & x $\to$ s3 & 1 & 1 & $-0.477$ & $-0.467$ [$-0.482$, $-0.453$] & [$-0.528$, $-0.440$] & [$-0.528$, $-0.440$] & [$-0.477$, $-0.477$] & never & never & -- \\
Thoemmes\_2013 & semi-s. & x $\to$ y & 1 & 1 & $-0.923$ & $-0.913$ [$-0.928$, $-0.898$] & [$-1.002$, $-0.871$] & [$-1.002$, $-0.871$] & [$-0.923$, $-0.923$] & never & never & -- \\
Thoemmes\_2013 & semi-s. & x $\to$ z & 1 & 1 & $\phantom{-}0.422$ & $\phantom{-}0.424$ [$\phantom{-}0.414$, $\phantom{-}0.434$] & [$\phantom{-}0.407$, $\phantom{-}0.448$] & [$\phantom{-}0.407$, $\phantom{-}0.448$] & [$\phantom{-}0.422$, $\phantom{-}0.422$] & never & never & -- \\
Thoemmes\_2013 & semi-s. & x $\to$ z2 & 1 & 1 & $\phantom{-}0.201$ & $\phantom{-}0.198$ [$\phantom{-}0.188$, $\phantom{-}0.208$] & [$\phantom{-}0.178$, $\phantom{-}0.219$] & [$\phantom{-}0.178$, $\phantom{-}0.219$] & [$\phantom{-}0.201$, $\phantom{-}0.201$] & never & never & -- \\
Thoemmes\_2013 & semi-s. & x $\to$ z3 & 1 & 1 & $-1.009$ & $-0.994$ [$-1.012$, $-0.975$] & [$-1.084$, $-0.947$] & [$-1.084$, $-0.947$] & [$-1.009$, $-1.009$] & never & never & -- \\
alarm & semi-s. & ANAPHYLAXIS $\to$ BP & 4 & 1 & $\phantom{-}0.108$ & $\phantom{-}0.096$ [$\phantom{-}0.076$, $\phantom{-}0.116$] & [$-0.068$, $\phantom{-}0.125$] & [$-0.068$, $\phantom{-}0.125$] & [$-0.000$, $\phantom{-}0.108$] & 1 & 1 & ANAPHYLAXIS$\to$TPR \\
alarm & semi-s. & ANAPHYLAXIS $\to$ HRBP & 4 & 1 & $-0.886$ & $-0.902$ [$-0.932$, $-0.872$] & [$-1.103$, $\phantom{-}0.122$] & [$-1.103$, $\phantom{-}0.122$] & [$-0.886$, $\phantom{-}0.000$] & 1 & 1 & ANAPHYLAXIS$\to$TPR \\
alarm & semi-s. & LVFAILURE $\to$ CO & 4 & 1 & $\phantom{-}1.421$ & $\phantom{-}1.447$ [$\phantom{-}1.427$, $\phantom{-}1.468$] & [$\phantom{-}1.322$, $\phantom{-}1.537$] & [$\phantom{-}1.322$, $\phantom{-}1.537$] & [$\phantom{-}1.421$, $\phantom{-}1.421$] & never & $>$3 & -- \\
alarm & semi-s. & LVFAILURE $\to$ PCWP & 4 & 1 & $-0.651$ & $-0.649$ [$-0.665$, $-0.632$] & [$-0.700$, $-0.629$] & [$-0.700$, $-0.629$] & [$-0.651$, $-0.651$] & never & $>$3 & -- \\
alarm & semi-s. & MINVOLSET $\to$ CATECHOL & 4 & 1 & $-0.608$ & $-0.591$ [$-0.629$, $-0.552$] & [$-0.675$, $\phantom{-}0.193$] & [$-0.675$, $\phantom{-}0.193$] & [$-0.608$, $-0.000$] & 1 & 1 & MINVOLSET$\to$VENTMACH \\
alarm & semi-s. & MINVOLSET $\to$ HRBP & 4 & 1 & $\phantom{-}0.835$ & $\phantom{-}0.804$ [$\phantom{-}0.749$, $\phantom{-}0.860$] & [$-0.270$, $\phantom{-}0.930$] & [$-0.270$, $\phantom{-}0.930$] & [$\phantom{-}0.000$, $\phantom{-}0.835$] & 1 & 1 & MINVOLSET$\to$VENTMACH \\
alarm & semi-s. & MINVOLSET $\to$ PRESS & 4 & 1 & $\phantom{-}1.360$ & $\phantom{-}1.355$ [$\phantom{-}1.326$, $\phantom{-}1.385$] & [$-0.056$, $\phantom{-}1.369$] & [$-0.056$, $\phantom{-}1.369$] & [$\phantom{-}0.000$, $\phantom{-}1.360$] & 1 & 1 & MINVOLSET$\to$VENTMACH \\
alarm & semi-s. & MINVOLSET $\to$ VENTLUNG & 4 & 1 & $\phantom{-}0.681$ & $\phantom{-}0.678$ [$\phantom{-}0.659$, $\phantom{-}0.697$] & [$-0.084$, $\phantom{-}0.702$] & [$-0.084$, $\phantom{-}0.702$] & [$\phantom{-}0.000$, $\phantom{-}0.681$] & 1 & 1 & MINVOLSET$\to$VENTMACH \\
alarm & semi-s. & PULMEMBOLUS $\to$ HR & 4 & 1 & $\phantom{-}2.287$ & $\phantom{-}2.281$ [$\phantom{-}2.229$, $\phantom{-}2.333$] & [$\phantom{-}2.070$, $\phantom{-}2.396$] & [$\phantom{-}2.070$, $\phantom{-}2.396$] & [$\phantom{-}2.287$, $\phantom{-}2.287$] & never & $>$3 & -- \\
alarm & semi-s. & PULMEMBOLUS $\to$ PAP & 4 & 1 & $\phantom{-}0.578$ & $\phantom{-}0.581$ [$\phantom{-}0.567$, $\phantom{-}0.594$] & [$\phantom{-}0.000$, $\phantom{-}0.594$] & [$\phantom{-}0.000$, $\phantom{-}0.594$] & [$\phantom{-}0.000$, $\phantom{-}0.578$] & 1 & 1 & PULMEMBOLUS$\to$PAP \\
alarm & semi-s. & TPR $\to$ BP & 4 & 1 & $-0.117$ & $-0.105$ [$-0.120$, $-0.091$] & [$-0.120$, $-0.004$] & [$-0.120$, $-0.004$] & [$-0.117$, $-0.117$] & never & $>$3 & -- \\
alarm & semi-s. & TPR $\to$ HRBP & 4 & 1 & $\phantom{-}0.952$ & $\phantom{-}0.977$ [$\phantom{-}0.958$, $\phantom{-}0.997$] & [$\phantom{-}0.889$, $\phantom{-}1.137$] & [$\phantom{-}0.889$, $\phantom{-}1.137$] & [$\phantom{-}0.952$, $\phantom{-}0.952$] & never & $>$3 & -- \\
alarm & semi-s. & VENTMACH $\to$ BP & 4 & 1 & $\phantom{-}0.295$ & $\phantom{-}0.292$ [$\phantom{-}0.274$, $\phantom{-}0.311$] & [$\phantom{-}0.273$, $\phantom{-}0.386$] & [$\phantom{-}0.273$, $\phantom{-}0.386$] & [$\phantom{-}0.295$, $\phantom{-}0.295$] & never & $>$3 & -- \\
alarm & semi-s. & VENTMACH $\to$ HR & 4 & 1 & $-0.781$ & $-0.753$ [$-0.791$, $-0.715$] & [$-0.977$, $-0.712$] & [$-0.977$, $-0.712$] & [$-0.781$, $-0.781$] & never & $>$3 & -- \\
alarm & semi-s. & VENTMACH $\to$ MINVOL & 4 & 1 & $-0.445$ & $-0.443$ [$-0.456$, $-0.431$] & [$-0.482$, $-0.430$] & [$-0.482$, $-0.430$] & [$-0.445$, $-0.445$] & never & $>$3 & -- \\
andes & semi-s. & GIVEN\_\allowbreak{}1 $\to$ SNode\_\allowbreak{}100 & 10 & 1 & $\phantom{-}0.187$ & $\phantom{-}0.189$ [$\phantom{-}0.137$, $\phantom{-}0.240$] & [$-0.749$, $\phantom{-}1.045$] & [$-0.749$, $\phantom{-}1.045$] & [$-0.000$, $\phantom{-}0.187$] & 1 & 1 & -- \\
andes & semi-s. & RApp2 $\to$ GOAL\_\allowbreak{}153 & 10 & 1 & $\phantom{-}0.218$ & $-0.026$ [$-0.376$, $\phantom{-}0.324$] & [$-1.053$, $\phantom{-}0.974$] & [$-1.053$, $\phantom{-}0.974$] & [$\phantom{-}0.218$, $\phantom{-}0.218$] & never & $>$3 & -- \\
andes & semi-s. & RApp2 $\to$ SNode\_\allowbreak{}43 & 10 & 1 & $\phantom{-}0.504$ & $\phantom{-}0.506$ [$\phantom{-}0.493$, $\phantom{-}0.519$] & [$\phantom{-}0.459$, $\phantom{-}0.595$] & [$\phantom{-}0.459$, $\phantom{-}0.595$] & [$\phantom{-}0.504$, $\phantom{-}0.504$] & never & $>$3 & -- \\
andes & semi-s. & TRY11 $\to$ GOAL\_\allowbreak{}81 & 10 & 4 & $-3.566$ & $-3.632$ [$-3.956$, $-3.307$] & [$-4.525$, $-2.636$] & [$-4.525$, $-2.636$] & [$-3.566$, $-3.566$] & $>$3 & $>$3 & -- \\
andes & semi-s. & TRY11 $\to$ SNode\_\allowbreak{}43 & 10 & 4 & $\phantom{-}0.011$ & $\phantom{-}0.009$ [$-0.010$, $\phantom{-}0.027$] & [$-0.084$, $\phantom{-}0.035$] & [$-0.084$, $\phantom{-}0.035$] & [$\phantom{-}0.011$, $\phantom{-}0.011$] & $>$3 & $>$3 & -- \\
andes & semi-s. & TRY12 $\to$ GOAL\_\allowbreak{}150 & 10 & 4 & $-43.982$ & $-45.462$ [$-50.673$, $-40.250$] & [$-62.203$, $\phantom{-}14.089$] & [$-62.203$, $\phantom{-}14.089$] & [$-43.982$, $-0.000$] & 1 & 1 & TRY12$\to$TRY11 \\
andes & semi-s. & TRY12 $\to$ SNode\_\allowbreak{}134 & 10 & 4 & $\phantom{-}10.027$ & $\phantom{-}10.267$ [$\phantom{-}9.155$, $\phantom{-}11.380$] & [$-2.760$, $\phantom{-}14.674$] & [$-2.760$, $\phantom{-}14.674$] & [$-0.000$, $\phantom{-}10.027$] & 1 & 1 & TRY12$\to$TRY11 \\
andes & semi-s. & TRY24 $\to$ GOAL\_\allowbreak{}113 & 10 & 2 & $-0.352$ & $-0.481$ [$-0.684$, $-0.278$] & [$-0.906$, $\phantom{-}0.808$] & [$-0.906$, $\phantom{-}0.808$] & [$-0.352$, $-0.352$] & never & $>$3 & -- \\
andes & semi-s. & TRY24 $\to$ SNode\_\allowbreak{}151 & 10 & 2 & $-7.307$ & $-8.164$ [$-9.509$, $-6.819$] & [$-14.574$, $\phantom{-}9.027$] & [$-14.574$, $\phantom{-}9.027$] & [$-7.307$, $-7.307$] & never & $>$3 & -- \\
andes & semi-s. & TRY25 $\to$ RApp8 & 10 & 2 & $-0.723$ & $-1.412$ [$-1.939$, $-0.884$] & [$-3.841$, $-0.422$] & [$-3.841$, $-0.422$] & [$-0.723$, $-0.000$] & 1 & 1 & -- \\
asia & semi-s. & asia $\to$ dysp & 3 & 1 & $-0.743$ & $-0.758$ [$-0.786$, $-0.730$] & [$-0.790$, $\phantom{-}0.068$] & [$-0.790$, $\phantom{-}0.068$] & [$-0.743$, $\phantom{-}0.000$] & 1 & 1 & asia$\to$tub \\
asia & semi-s. & asia $\to$ either & 3 & 1 & $\phantom{-}0.536$ & $\phantom{-}0.547$ [$\phantom{-}0.529$, $\phantom{-}0.565$] & [$-0.046$, $\phantom{-}0.565$] & [$-0.046$, $\phantom{-}0.565$] & [$-0.000$, $\phantom{-}0.536$] & 1 & 1 & asia$\to$tub \\
asia & semi-s. & asia $\to$ tub & 3 & 1 & $\phantom{-}0.667$ & $\phantom{-}0.676$ [$\phantom{-}0.662$, $\phantom{-}0.689$] & [$\phantom{-}0.000$, $\phantom{-}0.689$] & [$\phantom{-}0.000$, $\phantom{-}0.689$] & [$\phantom{-}0.000$, $\phantom{-}0.667$] & 1 & 1 & asia$\to$tub \\
asia & semi-s. & asia $\to$ xray & 3 & 1 & $\phantom{-}0.418$ & $\phantom{-}0.425$ [$\phantom{-}0.405$, $\phantom{-}0.445$] & [$-0.044$, $\phantom{-}0.442$] & [$-0.044$, $\phantom{-}0.442$] & [$-0.000$, $\phantom{-}0.418$] & 1 & 1 & asia$\to$tub \\
asia & semi-s. & bronc $\to$ dysp & 3 & 2 & $\phantom{-}1.065$ & $\phantom{-}1.063$ [$\phantom{-}1.052$, $\phantom{-}1.074$] & [$-0.142$, $\phantom{-}1.116$] & [$-0.142$, $\phantom{-}1.116$] & [$-0.138$, $\phantom{-}1.065$] & 1 & 1 & smoke$\to$bronc \\
asia & semi-s. & lung $\to$ dysp & 3 & 2 & $-1.809$ & $-1.794$ [$-1.811$, $-1.777$] & [$-1.844$, $-1.227$] & [$-1.844$, $-1.227$] & [$-1.809$, $-1.264$] & never & 1 & -- \\
asia & semi-s. & lung $\to$ either & 3 & 2 & $\phantom{-}1.305$ & $\phantom{-}1.300$ [$\phantom{-}1.291$, $\phantom{-}1.308$] & [$\phantom{-}1.281$, $\phantom{-}1.328$] & [$\phantom{-}1.281$, $\phantom{-}1.328$] & [$\phantom{-}1.305$, $\phantom{-}1.305$] & never & never & -- \\
asia & semi-s. & lung $\to$ xray & 3 & 2 & $\phantom{-}1.017$ & $\phantom{-}1.016$ [$\phantom{-}1.006$, $\phantom{-}1.027$] & [$\phantom{-}0.999$, $\phantom{-}1.045$] & [$\phantom{-}0.999$, $\phantom{-}1.045$] & [$\phantom{-}1.017$, $\phantom{-}1.017$] & never & never & -- \\
asia & semi-s. & smoke $\to$ bronc & 3 & 2 & $\phantom{-}1.067$ & $\phantom{-}1.058$ [$\phantom{-}1.044$, $\phantom{-}1.072$] & [$\phantom{-}0.000$, $\phantom{-}1.088$] & [$\phantom{-}0.000$, $\phantom{-}1.088$] & [$\phantom{-}0.000$, $\phantom{-}1.067$] & 1 & 1 & smoke$\to$bronc \\
asia & semi-s. & smoke $\to$ dysp & 3 & 2 & $-1.274$ & $-1.271$ [$-1.309$, $-1.234$] & [$-2.440$, $\phantom{-}1.226$] & [$-2.440$, $\phantom{-}1.226$] & [$-2.410$, $\phantom{-}1.136$] & 1 & 1 & smoke$\to$lung \\
asia & semi-s. & smoke $\to$ either & 3 & 2 & $\phantom{-}1.738$ & $\phantom{-}1.732$ [$\phantom{-}1.710$, $\phantom{-}1.755$] & [$-0.065$, $\phantom{-}1.761$] & [$-0.065$, $\phantom{-}1.761$] & [$-0.000$, $\phantom{-}1.738$] & 1 & 1 & smoke$\to$lung \\
asia & semi-s. & smoke $\to$ lung & 3 & 2 & $\phantom{-}1.332$ & $\phantom{-}1.334$ [$\phantom{-}1.320$, $\phantom{-}1.347$] & [$\phantom{-}0.000$, $\phantom{-}1.359$] & [$\phantom{-}0.000$, $\phantom{-}1.359$] & [$\phantom{-}0.000$, $\phantom{-}1.332$] & 1 & 1 & smoke$\to$lung \\
asia & semi-s. & smoke $\to$ xray & 3 & 2 & $\phantom{-}1.354$ & $\phantom{-}1.354$ [$\phantom{-}1.331$, $\phantom{-}1.376$] & [$-0.062$, $\phantom{-}1.385$] & [$-0.062$, $\phantom{-}1.385$] & [$-0.000$, $\phantom{-}1.354$] & 1 & 1 & smoke$\to$lung \\
barley & semi-s. & dg25 $\to$ aks\_\allowbreak{}m2 & 7 & 2 & $-1.188$ & $-1.159$ [$-1.185$, $-1.134$] & [$-1.671$, $-0.372$] & [$-1.671$, $-0.372$] & [$-1.188$, $-0.626$] & never & 1 & -- \\
barley & semi-s. & dg25 $\to$ s2225 & 7 & 2 & $-1.774$ & $-1.723$ [$-1.764$, $-1.683$] & [$-2.497$, $-0.555$] & [$-2.497$, $-0.555$] & [$-1.774$, $-0.950$] & never & 1 & -- \\
barley & semi-s. & frspdag $\to$ spndx & 7 & 2 & $-0.457$ & $-0.519$ [$-0.567$, $-0.470$] & [$-1.862$, $\phantom{-}0.773$] & [$-1.862$, $\phantom{-}0.773$] & [$-0.457$, $\phantom{-}0.056$] & 1 & 1 & -- \\
barley & semi-s. & jordtype $\to$ bgbyg & 7 & 1 & $\phantom{-}1.585$ & $\phantom{-}1.577$ [$\phantom{-}1.540$, $\phantom{-}1.613$] & [$-0.029$, $\phantom{-}1.618$] & [$-0.029$, $\phantom{-}1.618$] & [$\phantom{-}0.000$, $\phantom{-}1.585$] & 1 & 1 & jordtype$\to$rokap \\
barley & semi-s. & jordtype $\to$ nopt & 7 & 1 & $\phantom{-}0.931$ & $\phantom{-}0.946$ [$\phantom{-}0.921$, $\phantom{-}0.971$] & [$\phantom{-}0.910$, $\phantom{-}1.022$] & [$\phantom{-}0.910$, $\phantom{-}1.022$] & [$\phantom{-}0.931$, $\phantom{-}0.931$] & never & $>$3 & -- \\
barley & semi-s. & komm $\to$ aks\_\allowbreak{}m2 & 7 & 1 & $-8.425$ & $-8.406$ [$-8.558$, $-8.254$] & [$-8.951$, $-3.573$] & [$-8.951$, $-3.573$] & [$-8.425$, $-3.643$] & never & 1 & -- \\
barley & semi-s. & komm $\to$ ntilg & 7 & 1 & $\phantom{-}5.735$ & $\phantom{-}5.729$ [$\phantom{-}5.626$, $\phantom{-}5.832$] & [$\phantom{-}2.453$, $\phantom{-}6.110$] & [$\phantom{-}2.453$, $\phantom{-}6.110$] & [$\phantom{-}2.480$, $\phantom{-}5.735$] & never & 1 & -- \\
barley & semi-s. & nedbarea $\to$ slt22 & 7 & 1 & $\phantom{-}4.186$ & $\phantom{-}4.196$ [$\phantom{-}4.100$, $\phantom{-}4.291$] & [$\phantom{-}3.545$, $\phantom{-}6.231$] & [$\phantom{-}3.545$, $\phantom{-}6.231$] & [$\phantom{-}4.186$, $\phantom{-}5.986$] & never & 1 & -- \\
barley & semi-s. & rokap $\to$ keraks & 7 & 1 & $-3.015$ & $-3.031$ [$-3.061$, $-3.000$] & [$-10.350$, $-2.804$] & [$-10.350$, $-2.804$] & [$-9.983$, $-3.015$] & never & 1 & -- \\
barley & semi-s. & rokap $\to$ s2225 & 7 & 1 & $-3.018$ & $-3.022$ [$-3.053$, $-2.990$] & [$-16.911$, $-2.572$] & [$-16.911$, $-2.572$] & [$-16.480$, $-3.018$] & never & 1 & -- \\
barley & semi-s. & saatid $\to$ jordinf & 7 & 2 & $-0.382$ & $-0.378$ [$-0.393$, $-0.362$] & [$-0.399$, $\phantom{-}0.028$] & [$-0.399$, $\phantom{-}0.028$] & [$-0.382$, $\phantom{-}0.000$] & 1 & 1 & saatid$\to$frspdag \\
barley & semi-s. & sort $\to$ aks\_\allowbreak{}m2 & 7 & 3 & $-1.165$ & $-1.160$ [$-1.182$, $-1.139$] & [$-2.204$, $\phantom{-}0.290$] & [$-2.204$, $\phantom{-}0.290$] & [$-1.165$, $\phantom{-}0.000$] & 1 & 1 & sort$\to$sorttkv \\
barley & semi-s. & sort $\to$ udb & 7 & 3 & $-0.309$ & $-0.312$ [$-0.336$, $-0.288$] & [$-1.419$, $\phantom{-}0.292$] & [$-1.419$, $\phantom{-}0.292$] & [$-0.309$, $\phantom{-}0.000$] & 1 & 1 & sort$\to$sorttkv \\
barley & semi-s. & sorttkv $\to$ tkv & 7 & 3 & $\phantom{-}1.748$ & $\phantom{-}1.763$ [$\phantom{-}1.735$, $\phantom{-}1.790$] & [$\phantom{-}0.571$, $\phantom{-}2.846$] & [$\phantom{-}0.571$, $\phantom{-}2.846$] & [$\phantom{-}1.748$, $\phantom{-}1.748$] & never & $>$3 & -- \\
child & semi-s. & Birth\allowbreak{}Asphyxia $\to$ Cardiac\allowbreak{}Mixing & 2 & 2 & $\phantom{-}1.061$ & $\phantom{-}1.064$ [$\phantom{-}1.045$, $\phantom{-}1.083$] & [$-0.027$, $\phantom{-}1.083$] & [$-0.027$, $\phantom{-}1.083$] & [$\phantom{-}0.000$, $\phantom{-}1.061$] & 1 & 1 & BirthAsphyxia$\to$Disease \\
child & semi-s. & Birth\allowbreak{}Asphyxia $\to$ Lung\allowbreak{}Flow & 2 & 2 & $\phantom{-}1.461$ & $\phantom{-}1.463$ [$\phantom{-}1.440$, $\phantom{-}1.485$] & [$-0.015$, $\phantom{-}1.485$] & [$-0.015$, $\phantom{-}1.485$] & [$\phantom{-}0.000$, $\phantom{-}1.461$] & 1 & 1 & BirthAsphyxia$\to$Disease \\
child & semi-s. & Disease $\to$ Cardiac\allowbreak{}Mixing & 2 & 2 & $\phantom{-}0.948$ & $\phantom{-}0.947$ [$\phantom{-}0.938$, $\phantom{-}0.956$] & [$\phantom{-}0.000$, $\phantom{-}0.963$] & [$\phantom{-}0.000$, $\phantom{-}0.963$] & [$\phantom{-}0.000$, $\phantom{-}0.948$] & 1 & 1 & BirthAsphyxia$\to$Disease \\
child & semi-s. & Disease $\to$ Lung\allowbreak{}Parench & 2 & 2 & $\phantom{-}0.708$ & $\phantom{-}0.706$ [$\phantom{-}0.697$, $\phantom{-}0.715$] & [$\phantom{-}0.000$, $\phantom{-}0.748$] & [$\phantom{-}0.000$, $\phantom{-}0.748$] & [$\phantom{-}0.000$, $\phantom{-}0.708$] & 1 & 1 & BirthAsphyxia$\to$Disease \\
child & semi-s. & Lung\allowbreak{}Flow $\to$ Chest\allowbreak{}Xray & 2 & 2 & $\phantom{-}0.597$ & $\phantom{-}0.601$ [$\phantom{-}0.593$, $\phantom{-}0.609$] & [$\phantom{-}0.587$, $\phantom{-}1.012$] & [$\phantom{-}0.587$, $\phantom{-}1.012$] & [$\phantom{-}0.597$, $\phantom{-}1.000$] & never & 1 & -- \\
child & semi-s. & Lung\allowbreak{}Parench $\to$ Grunting & 2 & 2 & $\phantom{-}0.790$ & $\phantom{-}0.793$ [$\phantom{-}0.781$, $\phantom{-}0.805$] & [$\phantom{-}0.778$, $\phantom{-}1.607$] & [$\phantom{-}0.778$, $\phantom{-}1.607$] & [$\phantom{-}0.790$, $\phantom{-}1.583$] & never & 1 & -- \\
child & semi-s. & Sick $\to$ Age & 2 & 2 & $-1.103$ & $-1.111$ [$-1.125$, $-1.097$] & [$-1.736$, $\phantom{-}0.000$] & [$-1.736$, $\phantom{-}0.000$] & [$-1.723$, $\phantom{-}0.000$] & 1 & 1 & Sick$\to$Age \\
diabetes & semi-s. & bg\_\allowbreak{}0 $\to$ basal\_\allowbreak{}bal\_\allowbreak{}0 & 26 & 2 & $-2.412$ & $-2.428$ [$-2.477$, $-2.379$] & [$-2.495$, $-0.101$] & [$-2.495$, $-0.101$] & [$-2.412$, $-0.203$] & never & 1 & -- \\
diabetes & semi-s. & cho\_\allowbreak{}bal\_\allowbreak{}0 $\to$ ins\_\allowbreak{}dep\_\allowbreak{}10 & 26 & 1 & $\phantom{-}3.339$ & $\phantom{-}3.344$ [$\phantom{-}3.283$, $\phantom{-}3.406$] & [$\phantom{-}3.090$, $\phantom{-}3.841$] & [$\phantom{-}3.090$, $\phantom{-}3.841$] & [$\phantom{-}3.339$, $\phantom{-}3.671$] & never & 1 & -- \\
diabetes & semi-s. & cho\_\allowbreak{}bal\_\allowbreak{}12 $\to$ endo\_\allowbreak{}bal\_\allowbreak{}14 & 26 & 1 & $\phantom{-}3.030$ & $\phantom{-}3.057$ [$\phantom{-}3.001$, $\phantom{-}3.113$] & [$\phantom{-}2.541$, $\phantom{-}5.993$] & [$\phantom{-}2.541$, $\phantom{-}5.993$] & [$\phantom{-}3.030$, $\phantom{-}3.332$] & never & 1 & -- \\
diabetes & semi-s. & cho\_\allowbreak{}bal\_\allowbreak{}15 $\to$ glu\_\allowbreak{}prod\_\allowbreak{}22 & 26 & 1 & $-1.068$ & $-1.068$ [$-1.106$, $-1.029$] & [$-1.566$, $\phantom{-}1.660$] & [$-1.566$, $\phantom{-}1.660$] & [$-1.068$, $-0.213$] & never & 1 & -- \\
diabetes & semi-s. & cho\_\allowbreak{}bal\_\allowbreak{}2 $\to$ bg\_\allowbreak{}6 & 26 & 1 & $\phantom{-}3.273$ & $\phantom{-}3.284$ [$\phantom{-}3.244$, $\phantom{-}3.324$] & [$\phantom{-}2.693$, $\phantom{-}3.636$] & [$\phantom{-}2.693$, $\phantom{-}3.636$] & [$\phantom{-}2.791$, $\phantom{-}3.273$] & never & 1 & -- \\
diabetes & semi-s. & cho\_\allowbreak{}bal\_\allowbreak{}3 $\to$ cho\_\allowbreak{}bal\_\allowbreak{}17 & 26 & 1 & $\phantom{-}0.009$ & $\phantom{-}0.010$ [$-0.001$, $\phantom{-}0.020$] & [$-0.037$, $\phantom{-}0.139$] & [$-0.037$, $\phantom{-}0.139$] & [$\phantom{-}0.009$, $\phantom{-}0.009$] & never & $\geq$2 & -- \\
diabetes & semi-s. & cho\_\allowbreak{}bal\_\allowbreak{}4 $\to$ ins\_\allowbreak{}indep\_\allowbreak{}util\_\allowbreak{}8 & 26 & 1 & $-1.243$ & $-1.242$ [$-1.260$, $-1.224$] & [$-2.804$, $-0.449$] & [$-2.804$, $-0.449$] & [$-1.737$, $-1.243$] & never & 1 & -- \\
diabetes & semi-s. & cho\_\allowbreak{}bal\_\allowbreak{}6 $\to$ ins\_\allowbreak{}dep\_\allowbreak{}13 & 26 & 1 & $-42.307$ & $-42.167$ [$-42.479$, $-41.854$] & [$-45.240$, $-40.832$] & [$-45.240$, $-40.832$] & [$-42.307$, $-41.042$] & never & 1 & -- \\
diabetes & semi-s. & cho\_\allowbreak{}bal\_\allowbreak{}8 $\to$ gut\_\allowbreak{}abs\_\allowbreak{}10 & 26 & 1 & $\phantom{-}3.356$ & $\phantom{-}3.367$ [$\phantom{-}3.353$, $\phantom{-}3.382$] & [$\phantom{-}3.310$, $\phantom{-}3.434$] & [$\phantom{-}3.310$, $\phantom{-}3.434$] & [$\phantom{-}3.356$, $\phantom{-}3.356$] & never & $\geq$2 & -- \\
diabetes & semi-s. & gut\_\allowbreak{}abs\_\allowbreak{}0 $\to$ endo\_\allowbreak{}bal\_\allowbreak{}3 & 26 & 1 & $-6.742$ & $-6.923$ [$-7.134$, $-6.712$] & [$-7.103$, $-0.862$] & [$-7.103$, $-0.862$] & [$-6.742$, $-1.214$] & never & 1 & -- \\
diabetes & semi-s. & gut\_\allowbreak{}abs\_\allowbreak{}1 $\to$ ins\_\allowbreak{}indep\_\allowbreak{}21 & 26 & 1 & $\phantom{-}4.276$ & $\phantom{-}4.225$ [$\phantom{-}3.929$, $\phantom{-}4.521$] & [$\phantom{-}2.411$, $\phantom{-}4.293$] & [$\phantom{-}2.411$, $\phantom{-}4.293$] & [$\phantom{-}3.399$, $\phantom{-}4.276$] & never & 1 & -- \\
diabetes & semi-s. & gut\_\allowbreak{}abs\_\allowbreak{}11 $\to$ tot\_\allowbreak{}bal\_\allowbreak{}18 & 26 & 1 & $-0.465$ & $-0.645$ [$-1.024$, $-0.266$] & [$-5.184$, $\phantom{-}0.966$] & [$-5.184$, $\phantom{-}0.966$] & [$-0.465$, $\phantom{-}0.643$] & 1 & 1 & gut\_abs\_11$\to$cho\_bal\_11 \\
diabetes & semi-s. & gut\_\allowbreak{}abs\_\allowbreak{}14 $\to$ tot\_\allowbreak{}bal\_\allowbreak{}19 & 26 & 1 & $-4.780$ & $-5.317$ [$-5.880$, $-4.754$] & [$-9.102$, $\phantom{-}10.428$] & [$-9.102$, $\phantom{-}10.428$] & [$-4.780$, $\phantom{-}1.191$] & 1 & 1 & gut\_abs\_14$\to$cho\_bal\_14 \\
diabetes & semi-s. & gut\_\allowbreak{}abs\_\allowbreak{}19 $\to$ ins\_\allowbreak{}indep\_\allowbreak{}23 & 26 & 1 & $\phantom{-}30.328$ & $\phantom{-}30.762$ [$\phantom{-}30.132$, $\phantom{-}31.392$] & [$-9.521$, $\phantom{-}33.208$] & [$-9.521$, $\phantom{-}33.208$] & [$\phantom{-}2.660$, $\phantom{-}30.328$] & never & 1 & gut\_abs\_19$\to$cho\_bal\_19 \\
diabetes & semi-s. & gut\_\allowbreak{}abs\_\allowbreak{}3 $\to$ basal\_\allowbreak{}bal\_\allowbreak{}5 & 26 & 1 & $\phantom{-}1.012$ & $\phantom{-}0.972$ [$\phantom{-}0.861$, $\phantom{-}1.084$] & [$-5.581$, $\phantom{-}1.701$] & [$-5.581$, $\phantom{-}1.701$] & [$-3.011$, $\phantom{-}1.012$] & 1 & 1 & -- \\
diabetes & semi-s. & gut\_\allowbreak{}abs\_\allowbreak{}4 $\to$ ins\_\allowbreak{}dep\_\allowbreak{}util\_\allowbreak{}11 & 26 & 1 & $-0.493$ & $-0.364$ [$-0.879$, $\phantom{-}0.151$] & [$-0.835$, $\phantom{-}2.065$] & [$-0.835$, $\phantom{-}2.065$] & [$-0.493$, $\phantom{-}0.498$] & 1 & 1 & -- \\
diabetes & semi-s. & gut\_\allowbreak{}abs\_\allowbreak{}6 $\to$ endo\_\allowbreak{}bal\_\allowbreak{}10 & 26 & 1 & $\phantom{-}19.736$ & $\phantom{-}19.386$ [$\phantom{-}18.355$, $\phantom{-}20.418$] & [$-6.422$, $\phantom{-}22.114$] & [$-6.422$, $\phantom{-}22.114$] & [$-3.960$, $\phantom{-}19.736$] & 1 & 1 & gut\_abs\_6$\to$cho\_bal\_6 \\
diabetes & semi-s. & gut\_\allowbreak{}abs\_\allowbreak{}8 $\to$ cho\_\allowbreak{}17 & 26 & 1 & $-0.013$ & $-0.013$ [$-0.040$, $\phantom{-}0.014$] & [$-0.032$, $\phantom{-}0.055$] & [$-0.032$, $\phantom{-}0.055$] & [$-0.013$, $\phantom{-}0.000$] & 1 & 1 & -- \\
diabetes & semi-s. & ins\_\allowbreak{}indep\_\allowbreak{}util\_\allowbreak{}0 $\to$ endo\_\allowbreak{}bal\_\allowbreak{}17 & 26 & 2 & $-28.237$ & $-28.389$ [$-30.667$, $-26.111$] & [$-59.114$, $-23.979$] & [$-59.114$, $-23.979$] & [$-53.027$, $-28.237$] & never & 1 & -- \\
hailfinder & semi-s. & Date $\to$ AMCINIn\allowbreak{}Scen & 1 & 1 & $\phantom{-}1.407$ & $\phantom{-}1.400$ [$\phantom{-}1.377$, $\phantom{-}1.422$] & [$-0.024$, $\phantom{-}1.424$] & [$-0.024$, $\phantom{-}1.424$] & [$\phantom{-}0.000$, $\phantom{-}1.407$] & 1 & 1 & Date$\to$Scenario \\
hailfinder & semi-s. & Date $\to$ Wind\allowbreak{}Aloft & 1 & 1 & $-1.607$ & $-1.614$ [$-1.636$, $-1.593$] & [$-1.636$, $\phantom{-}0.018$] & [$-1.636$, $\phantom{-}0.018$] & [$-1.607$, $-0.000$] & 1 & 1 & Date$\to$Scenario \\
hailfinder & semi-s. & Scenario $\to$ R5Fcst & 1 & 1 & $-0.275$ & $-0.261$ [$-0.276$, $-0.247$] & [$-0.888$, $\phantom{-}0.905$] & [$-0.888$, $\phantom{-}0.905$] & [$-0.817$, $\phantom{-}0.807$] & 1 & 1 & Date$\to$Scenario \\
hepar2 & semi-s. & choledocholithotomy $\to$ consciousness & 9 & 5 & $\phantom{-}1.911$ & $\phantom{-}1.916$ [$\phantom{-}1.898$, $\phantom{-}1.933$] & [$\phantom{-}1.829$, $\phantom{-}1.948$] & [$\phantom{-}1.829$, $\phantom{-}1.948$] & [$\phantom{-}1.911$, $\phantom{-}1.911$] & $>$3 & $>$3 & -- \\
hepar2 & semi-s. & choledocholithotomy $\to$ irregular\_\allowbreak{}liver & 9 & 5 & $-4.500$ & $-4.520$ [$-4.554$, $-4.485$] & [$-4.587$, $-4.366$] & [$-4.587$, $-4.366$] & [$-4.500$, $-4.500$] & $>$3 & $>$3 & -- \\
hepar2 & semi-s. & diabetes $\to$ alcohol & 9 & 1 & $\phantom{-}0.584$ & $\phantom{-}0.579$ [$\phantom{-}0.559$, $\phantom{-}0.599$] & [$-0.077$, $\phantom{-}0.669$] & [$-0.077$, $\phantom{-}0.669$] & [$-0.000$, $\phantom{-}0.584$] & 1 & 1 & diabetes$\to$obesity \\
hepar2 & semi-s. & diabetes $\to$ itching & 9 & 1 & $-0.890$ & $-0.895$ [$-0.925$, $-0.865$] & [$-0.997$, $\phantom{-}0.087$] & [$-0.997$, $\phantom{-}0.087$] & [$-0.890$, $\phantom{-}0.000$] & 1 & 1 & diabetes$\to$obesity \\
hepar2 & semi-s. & gallstones $\to$ alt & 9 & 5 & $-0.382$ & $-0.408$ [$-0.431$, $-0.385$] & [$-0.601$, $\phantom{-}0.070$] & [$-0.601$, $\phantom{-}0.070$] & [$-0.382$, $-0.000$] & 1 & 1 & gallstones$\to$choledocholithotomy \\
hepar2 & semi-s. & gallstones $\to$ ggtp & 9 & 5 & $\phantom{-}2.868$ & $\phantom{-}2.887$ [$\phantom{-}2.849$, $\phantom{-}2.926$] & [$-0.039$, $\phantom{-}3.055$] & [$-0.039$, $\phantom{-}3.055$] & [$-0.000$, $\phantom{-}2.868$] & 1 & 1 & gallstones$\to$choledocholithotomy \\
hepar2 & semi-s. & gallstones $\to$ spleen & 9 & 5 & $\phantom{-}3.776$ & $\phantom{-}3.800$ [$\phantom{-}3.749$, $\phantom{-}3.851$] & [$-0.017$, $\phantom{-}3.971$] & [$-0.017$, $\phantom{-}3.971$] & [$\phantom{-}0.000$, $\phantom{-}3.776$] & 1 & 1 & gallstones$\to$choledocholithotomy \\
hepar2 & semi-s. & hepatotoxic $\to$ phosphatase & 9 & 1 & $-2.367$ & $-2.368$ [$-2.398$, $-2.339$] & [$-2.418$, $-0.969$] & [$-2.418$, $-0.969$] & [$-2.367$, $-1.015$] & never & 1 & -- \\
hepar2 & semi-s. & obesity $\to$ edema & 9 & 1 & $-0.486$ & $-0.482$ [$-0.494$, $-0.470$] & [$-0.542$, $-0.410$] & [$-0.542$, $-0.410$] & [$-0.486$, $-0.486$] & never & $>$3 & -- \\
hepar2 & semi-s. & obesity $\to$ triglycerides & 9 & 1 & $-0.952$ & $-0.942$ [$-0.957$, $-0.926$] & [$-0.965$, $-0.893$] & [$-0.965$, $-0.893$] & [$-0.952$, $-0.952$] & never & $>$3 & -- \\
insurance & semi-s. & Age $\to$ Accident & 3 & 3 & $-1.139$ & $-1.126$ [$-1.169$, $-1.084$] & [$-2.384$, $\phantom{-}1.514$] & [$-2.384$, $\phantom{-}1.514$] & [$-2.332$, $\phantom{-}1.448$] & 1 & 1 & Age$\to$SocioEcon \\
insurance & semi-s. & Age $\to$ Med\allowbreak{}Cost & 3 & 3 & $-1.809$ & $-1.781$ [$-1.912$, $-1.651$] & [$-5.303$, $\phantom{-}4.682$] & [$-5.303$, $\phantom{-}4.682$] & [$-5.159$, $\phantom{-}4.599$] & 1 & 1 & Age$\to$SocioEcon \\
insurance & semi-s. & Driving\allowbreak{}Skill $\to$ Other\allowbreak{}Car\allowbreak{}Cost & 3 & 3 & $-0.759$ & $-0.760$ [$-0.771$, $-0.749$] & [$-0.924$, $\phantom{-}0.702$] & [$-0.924$, $\phantom{-}0.702$] & [$-0.880$, $\phantom{-}0.667$] & 1 & 1 & SocioEcon$\to$RiskAversion \\
insurance & semi-s. & Make\allowbreak{}Model $\to$ Car\allowbreak{}Value & 3 & 3 & $-1.248$ & $-1.247$ [$-1.255$, $-1.238$] & [$-2.241$, $-1.204$] & [$-2.241$, $-0.027$] & [$-2.206$, $-0.028$] & never & 1 & -- \\
insurance & semi-s. & Make\allowbreak{}Model $\to$ Rugged\allowbreak{}Auto & 3 & 3 & $-1.295$ & $-1.294$ [$-1.302$, $-1.285$] & [$-2.256$, $-1.250$] & [$-2.256$, $-0.108$] & [$-2.229$, $-0.107$] & never & 1 & -- \\
insurance & semi-s. & Risk\allowbreak{}Aversion $\to$ Driv\allowbreak{}Hist & 3 & 3 & $\phantom{-}0.701$ & $\phantom{-}0.699$ [$\phantom{-}0.680$, $\phantom{-}0.718$] & [$\phantom{-}0.184$, $\phantom{-}2.223$] & [$-0.398$, $\phantom{-}2.223$] & [$-0.366$, $\phantom{-}2.201$] & 2 & 1 & Age$\to$SocioEcon \,/\, SocioEcon$\to$RiskAversion \\
insurance & semi-s. & Senior\allowbreak{}Train $\to$ Driv\allowbreak{}Quality & 3 & 3 & $-1.200$ & $-1.188$ [$-1.211$, $-1.165$] & [$-1.979$, $\phantom{-}0.476$] & [$-1.979$, $\phantom{-}0.476$] & [$-1.967$, $\phantom{-}0.468$] & 1 & 1 & SocioEcon$\to$RiskAversion \\
insurance & semi-s. & Socio\allowbreak{}Econ $\to$ Driving\allowbreak{}Skill & 3 & 3 & $-0.595$ & $-0.574$ [$-0.595$, $-0.552$] & [$-0.620$, $\phantom{-}1.805$] & [$-0.970$, $\phantom{-}1.805$] & [$-0.957$, $\phantom{-}1.732$] & 1 & 1 & Age$\to$SocioEcon \\
insurance & semi-s. & Socio\allowbreak{}Econ $\to$ This\allowbreak{}Car\allowbreak{}Cost & 3 & 3 & $\phantom{-}0.468$ & $\phantom{-}0.522$ [$\phantom{-}0.441$, $\phantom{-}0.602$] & [$-2.435$, $\phantom{-}1.411$] & [$-2.435$, $\phantom{-}2.463$] & [$-2.376$, $\phantom{-}2.382$] & 1 & 1 & Age$\to$SocioEcon \\
mediator & semi-s. & I $\to$ Y & 3 & 3 & $\phantom{-}1.360$ & $\phantom{-}1.356$ [$\phantom{-}1.344$, $\phantom{-}1.369$] & [$\phantom{-}1.343$, $\phantom{-}1.928$] & [$\phantom{-}0.000$, $\phantom{-}2.170$] & [$\phantom{-}0.000$, $\phantom{-}2.161$] & 2 & 1 & Z$\to$I \,/\, Z$\to$X \\
mediator & semi-s. & X $\to$ I & 3 & 3 & $-1.087$ & $-1.083$ [$-1.097$, $-1.069$] & [$-1.097$, $\phantom{-}0.000$] & [$-1.097$, $\phantom{-}0.342$] & [$-1.087$, $\phantom{-}0.324$] & 1 & 1 & X$\to$I \\
mediator & semi-s. & X $\to$ Y & 3 & 3 & $-2.603$ & $-2.600$ [$-2.623$, $-2.576$] & [$-2.623$, $-1.109$] & [$-2.623$, $\phantom{-}0.000$] & [$-2.603$, $\phantom{-}0.000$] & 2 & 1 & Z$\to$I \,/\, Z$\to$X \\
mediator & semi-s. & Z $\to$ I & 3 & 3 & $\phantom{-}0.642$ & $\phantom{-}0.632$ [$\phantom{-}0.612$, $\phantom{-}0.653$] & [$-0.672$, $\phantom{-}0.653$] & [$-0.672$, $\phantom{-}0.653$] & [$-0.645$, $\phantom{-}0.642$] & 1 & 1 & Z$\to$X \\
mediator & semi-s. & Z $\to$ X & 3 & 3 & $-1.185$ & $-1.184$ [$-1.197$, $-1.170$] & [$-1.197$, $\phantom{-}0.000$] & [$-1.197$, $\phantom{-}0.000$] & [$-1.185$, $\phantom{-}0.000$] & 1 & 1 & Z$\to$X \\
mediator & semi-s. & Z $\to$ Y & 3 & 3 & $\phantom{-}2.206$ & $\phantom{-}2.197$ [$\phantom{-}2.154$, $\phantom{-}2.240$] & [$-0.917$, $\phantom{-}2.240$] & [$-0.917$, $\phantom{-}2.240$] & [$-0.877$, $\phantom{-}2.206$] & 1 & 1 & Z$\to$X \\
munin1 & semi-s. & R\_\allowbreak{}LNLBE\_\allowbreak{}MED\_\allowbreak{}PATHO $\to$ R\_\allowbreak{}LNLBE\_\allowbreak{}APB\_\allowbreak{}DENERV & 6 & 4 & $-1.235$ & $-1.255$ [$-1.270$, $-1.241$] & [$-1.307$, $-1.208$] & [$-1.307$, $-1.208$] & [$-1.235$, $-1.235$] & $>$3 & $>$3 & -- \\
munin1 & semi-s. & R\_\allowbreak{}LNLBE\_\allowbreak{}MED\_\allowbreak{}PATHO $\to$ R\_\allowbreak{}MED\_\allowbreak{}AMP\_\allowbreak{}WA & 6 & 4 & $-2.329$ & $-2.308$ [$-2.370$, $-2.246$] & [$-2.619$, $-2.034$] & [$-2.619$, $-2.034$] & [$-2.329$, $-2.329$] & $>$3 & $>$3 & -- \\
munin1 & semi-s. & R\_\allowbreak{}LNLW\_\allowbreak{}MED\_\allowbreak{}PATHO $\to$ R\_\allowbreak{}APB\_\allowbreak{}SPONT\_\allowbreak{}INS\_\allowbreak{}ACT & 6 & 2 & $-1.695$ & $-1.689$ [$-1.762$, $-1.615$] & [$-2.108$, $-1.405$] & [$-2.108$, $-1.405$] & [$-1.695$, $-1.695$] & never & $>$3 & -- \\
munin1 & semi-s. & R\_\allowbreak{}LNLW\_\allowbreak{}MED\_\allowbreak{}PATHO $\to$ R\_\allowbreak{}MED\_\allowbreak{}ALLDEL\_\allowbreak{}WA & 6 & 2 & $\phantom{-}2.990$ & $\phantom{-}2.998$ [$\phantom{-}2.962$, $\phantom{-}3.033$] & [$\phantom{-}1.528$, $\phantom{-}3.107$] & [$\phantom{-}1.528$, $\phantom{-}3.107$] & [$\phantom{-}1.619$, $\phantom{-}2.990$] & never & 1 & -- \\
munin1 & semi-s. & R\_\allowbreak{}MEDD2\_\allowbreak{}RD\_\allowbreak{}WD $\to$ R\_\allowbreak{}MEDD2\_\allowbreak{}LSLOW\_\allowbreak{}WD & 6 & 2 & $-1.330$ & $-1.328$ [$-1.335$, $-1.321$] & [$-1.345$, $-1.177$] & [$-1.345$, $-1.177$] & [$-1.330$, $-1.203$] & never & 1 & -- \\
paths & semi-s. & 10 $\to$ 4 & 1 & 1 & $\phantom{-}0.591$ & $\phantom{-}0.585$ [$\phantom{-}0.567$, $\phantom{-}0.603$] & [$-0.038$, $\phantom{-}0.615$] & [$-0.038$, $\phantom{-}0.615$] & [$-0.000$, $\phantom{-}0.591$] & 1 & 1 & 15$\to$14 \\
paths & semi-s. & 12 $\to$ 1 & 1 & 1 & $\phantom{-}0.715$ & $\phantom{-}0.717$ [$\phantom{-}0.669$, $\phantom{-}0.766$] & [$-0.064$, $\phantom{-}0.802$] & [$-0.064$, $\phantom{-}0.802$] & [$\phantom{-}0.000$, $\phantom{-}0.715$] & 1 & 1 & 15$\to$14 \\
paths & semi-s. & 12 $\to$ 8 & 1 & 1 & $\phantom{-}0.353$ & $\phantom{-}0.351$ [$\phantom{-}0.338$, $\phantom{-}0.364$] & [$-0.024$, $\phantom{-}0.364$] & [$-0.024$, $\phantom{-}0.364$] & [$\phantom{-}0.000$, $\phantom{-}0.353$] & 1 & 1 & 15$\to$14 \\
paths & semi-s. & 13 $\to$ 4 & 1 & 1 & $-0.319$ & $-0.321$ [$-0.347$, $-0.295$] & [$-0.367$, $\phantom{-}0.034$] & [$-0.367$, $\phantom{-}0.034$] & [$-0.319$, $\phantom{-}0.000$] & 1 & 1 & 15$\to$14 \\
paths & semi-s. & 14 $\to$ 13 & 1 & 1 & $\phantom{-}0.801$ & $\phantom{-}0.802$ [$\phantom{-}0.791$, $\phantom{-}0.814$] & [$\phantom{-}0.000$, $\phantom{-}0.816$] & [$\phantom{-}0.000$, $\phantom{-}0.816$] & [$\phantom{-}0.000$, $\phantom{-}0.801$] & 1 & 1 & 15$\to$14 \\
paths & semi-s. & 15 $\to$ 1 & 1 & 1 & $-0.354$ & $-0.304$ [$-0.381$, $-0.226$] & [$-0.381$, $\phantom{-}0.147$] & [$-0.381$, $\phantom{-}0.147$] & [$-0.354$, $\phantom{-}0.000$] & 1 & 1 & 15$\to$14 \\
paths & semi-s. & 15 $\to$ 8 & 1 & 1 & $-0.175$ & $-0.168$ [$-0.190$, $-0.146$] & [$-0.190$, $\phantom{-}0.036$] & [$-0.190$, $\phantom{-}0.036$] & [$-0.175$, $\phantom{-}0.000$] & 1 & 1 & 15$\to$14 \\
paths & semi-s. & 4 $\to$ 1 & 1 & 1 & $\phantom{-}1.929$ & $\phantom{-}1.916$ [$\phantom{-}1.905$, $\phantom{-}1.927$] & [$-0.027$, $\phantom{-}1.941$] & [$-0.027$, $\phantom{-}1.941$] & [$-0.000$, $\phantom{-}1.929$] & 1 & 1 & 15$\to$14 \\
paths & semi-s. & 5 $\to$ 3 & 1 & 1 & $\phantom{-}1.831$ & $\phantom{-}1.824$ [$\phantom{-}1.812$, $\phantom{-}1.836$] & [$-0.015$, $\phantom{-}1.855$] & [$-0.015$, $\phantom{-}1.855$] & [$-0.000$, $\phantom{-}1.831$] & 1 & 1 & 15$\to$14 \\
paths & semi-s. & 7 $\to$ 1 & 1 & 1 & $-2.080$ & $-2.044$ [$-2.075$, $-2.013$] & [$-2.090$, $\phantom{-}0.066$] & [$-2.090$, $\phantom{-}0.066$] & [$-2.080$, $-0.000$] & 1 & 1 & 15$\to$14 \\
paths & semi-s. & 8 $\to$ 3 & 1 & 1 & $-1.573$ & $-1.561$ [$-1.590$, $-1.532$] & [$-1.596$, $\phantom{-}0.027$] & [$-1.596$, $\phantom{-}0.027$] & [$-1.573$, $\phantom{-}0.000$] & 1 & 1 & 15$\to$14 \\
sachs & semi-s. & Mek $\to$ Akt & 8 & 6 & $\phantom{-}1.229$ & $\phantom{-}1.223$ [$\phantom{-}1.210$, $\phantom{-}1.236$] & [$\phantom{-}1.196$, $\phantom{-}1.738$] & [$\phantom{-}1.196$, $\phantom{-}1.779$] & [$-0.047$, $\phantom{-}1.773$] & 3 & 1 & PKC$\to$Mek \,/\, PKC$\to$PKA \,/\, Raf$\to$Mek \\
sachs & semi-s. & PKA $\to$ Akt & 8 & 6 & $-5.107$ & $-5.129$ [$-5.159$, $-5.098$] & [$-5.175$, $-1.651$] & [$-5.175$, $-1.651$] & [$-5.107$, $\phantom{-}0.000$] & 3 & 1 & PKC$\to$Mek \,/\, PKC$\to$PKA \,/\, Raf$\to$Mek \\
sachs & semi-s. & PKA $\to$ Mek & 8 & 6 & $-2.800$ & $-2.813$ [$-2.832$, $-2.793$] & [$-2.845$, $\phantom{-}0.000$] & [$-2.845$, $\phantom{-}0.000$] & [$-2.800$, $\phantom{-}0.135$] & 1 & 1 & PKA$\to$Raf \\
sachs & semi-s. & PKA $\to$ Raf & 8 & 6 & $\phantom{-}1.315$ & $\phantom{-}1.323$ [$\phantom{-}1.309$, $\phantom{-}1.336$] & [$-0.101$, $\phantom{-}1.692$] & [$-0.101$, $\phantom{-}1.857$] & [$-0.085$, $\phantom{-}1.839$] & 1 & 1 & PKA$\to$Raf \\
sachs & semi-s. & PKC $\to$ Jnk & 8 & 6 & $-1.315$ & $-1.305$ [$-1.325$, $-1.285$] & [$-1.472$, $\phantom{-}0.000$] & [$-1.472$, $\phantom{-}0.000$] & [$-1.462$, $\phantom{-}0.000$] & 1 & 1 & PKA$\to$Jnk \\
sachs & semi-s. & PKC $\to$ PKA & 8 & 6 & $-0.755$ & $-0.749$ [$-0.763$, $-0.735$] & [$-0.900$, $\phantom{-}0.093$] & [$-0.900$, $\phantom{-}0.093$] & [$-0.894$, $\phantom{-}0.084$] & 1 & 1 & PKA$\to$Raf \\
sachs & semi-s. & Plcg $\to$ PIP3 & 8 & 2 & $\phantom{-}0.587$ & $\phantom{-}0.586$ [$\phantom{-}0.572$, $\phantom{-}0.600$] & [$-0.309$, $\phantom{-}0.600$] & [$-0.309$, $\phantom{-}0.600$] & [$-0.287$, $\phantom{-}0.587$] & 1 & 1 & PIP3$\to$PIP2 \\
water & semi-s. & CKNI\_\allowbreak{}12\_\allowbreak{}00 $\to$ CBODD\_\allowbreak{}12\_\allowbreak{}15 & 2 & 1 & $\phantom{-}0.672$ & $\phantom{-}0.668$ [$\phantom{-}0.654$, $\phantom{-}0.682$] & [$\phantom{-}0.581$, $\phantom{-}0.677$] & [$\phantom{-}0.581$, $\phantom{-}0.677$] & [$\phantom{-}0.672$, $\phantom{-}0.672$] & never & never & -- \\
water & semi-s. & CKNI\_\allowbreak{}12\_\allowbreak{}00 $\to$ CKND\_\allowbreak{}12\_\allowbreak{}15 & 2 & 1 & $\phantom{-}1.372$ & $\phantom{-}1.369$ [$\phantom{-}1.355$, $\phantom{-}1.383$] & [$\phantom{-}1.337$, $\phantom{-}1.412$] & [$\phantom{-}1.337$, $\phantom{-}1.412$] & [$\phantom{-}1.372$, $\phantom{-}1.372$] & never & never & -- \\
water & semi-s. & CKNI\_\allowbreak{}12\_\allowbreak{}00 $\to$ CKNI\_\allowbreak{}12\_\allowbreak{}45 & 2 & 1 & $-0.814$ & $-0.811$ [$-0.834$, $-0.787$] & [$-0.834$, $\phantom{-}0.026$] & [$-0.834$, $\phantom{-}0.026$] & [$-0.814$, $\phantom{-}0.000$] & 1 & 1 & CKNI\_12\_00$\to$CKNI\_12\_15 \\
water & semi-s. & CKNI\_\allowbreak{}12\_\allowbreak{}00 $\to$ CNON\_\allowbreak{}12\_\allowbreak{}45 & 2 & 1 & $-1.385$ & $-1.387$ [$-1.417$, $-1.356$] & [$-1.691$, $-1.325$] & [$-1.691$, $-1.325$] & [$-1.385$, $-1.385$] & never & never & -- \\
water & semi-s. & CKNI\_\allowbreak{}12\_\allowbreak{}15 $\to$ CBODN\_\allowbreak{}12\_\allowbreak{}45 & 2 & 1 & $-1.200$ & $-1.194$ [$-1.207$, $-1.181$] & [$-1.269$, $-0.995$] & [$-1.269$, $-0.995$] & [$-1.211$, $-1.200$] & never & 1 & -- \\
water & semi-s. & CKNI\_\allowbreak{}12\_\allowbreak{}15 $\to$ CKNI\_\allowbreak{}12\_\allowbreak{}30 & 2 & 1 & $-0.779$ & $-0.772$ [$-0.782$, $-0.762$] & [$-0.782$, $\phantom{-}0.000$] & [$-0.782$, $\phantom{-}0.000$] & [$-0.779$, $\phantom{-}0.000$] & 1 & 1 & CKNI\_12\_00$\to$CKNI\_12\_15 \\
water & semi-s. & CKNI\_\allowbreak{}12\_\allowbreak{}15 $\to$ CNOD\_\allowbreak{}12\_\allowbreak{}45 & 2 & 1 & $-1.671$ & $-1.672$ [$-1.688$, $-1.656$] & [$-1.807$, $-1.587$] & [$-1.807$, $-1.587$] & [$-1.743$, $-1.671$] & never & 1 & -- \\
water & semi-s. & C\_\allowbreak{}NI\_\allowbreak{}12\_\allowbreak{}00 $\to$ CBODD\_\allowbreak{}12\_\allowbreak{}30 & 2 & 1 & $-1.391$ & $-1.397$ [$-1.418$, $-1.375$] & [$-1.429$, $-0.739$] & [$-1.429$, $-0.739$] & [$-1.391$, $-0.831$] & never & 1 & -- \\
water & semi-s. & C\_\allowbreak{}NI\_\allowbreak{}12\_\allowbreak{}00 $\to$ CNOD\_\allowbreak{}12\_\allowbreak{}45 & 2 & 1 & $-0.885$ & $-0.880$ [$-0.908$, $-0.851$] & [$-0.868$, $-0.016$] & [$-0.868$, $-0.016$] & [$-0.885$, $-0.226$] & never & 1 & -- \\
water & semi-s. & C\_\allowbreak{}NI\_\allowbreak{}12\_\allowbreak{}15 $\to$ C\_\allowbreak{}NI\_\allowbreak{}12\_\allowbreak{}45 & 2 & 1 & $-0.662$ & $-0.659$ [$-0.674$, $-0.645$] & [$-0.674$, $\phantom{-}0.024$] & [$-0.674$, $\phantom{-}0.024$] & [$-0.662$, $-0.000$] & 1 & 1 & C\_NI\_12\_00$\to$C\_NI\_12\_15 \\
win95pts & semi-s. & App\allowbreak{}Dt\allowbreak{}Gn\allowbreak{}Tm $\to$ Desk\allowbreak{}Prnt\allowbreak{}Spd & 9 & 2 & $-1.226$ & $-1.230$ [$-1.240$, $-1.220$] & [$-1.517$, $-1.214$] & [$-1.517$, $-1.214$] & [$-1.485$, $-1.226$] & never & 1 & -- \\
win95pts & semi-s. & Avlbl\allowbreak{}Vrtl\allowbreak{}Mmry $\to$ Incmplt\allowbreak{}PS & 9 & 3 & $-1.148$ & $-1.144$ [$-1.156$, $-1.132$] & [$-1.174$, $-1.067$] & [$-1.176$, $-1.067$] & [$-1.148$, $-1.148$] & never & $>$3 & -- \\
win95pts & semi-s. & Prt\allowbreak{}Mem $\to$ Desk\allowbreak{}Prnt\allowbreak{}Spd & 9 & 1 & $\phantom{-}1.022$ & $\phantom{-}1.018$ [$\phantom{-}1.004$, $\phantom{-}1.032$] & [$\phantom{-}0.963$, $\phantom{-}1.092$] & [$\phantom{-}0.963$, $\phantom{-}1.092$] & [$\phantom{-}1.022$, $\phantom{-}1.022$] & never & $>$3 & -- \\
win95pts & semi-s. & Prt\allowbreak{}Mem $\to$ Nt\allowbreak{}Grbld & 9 & 1 & $-1.115$ & $-1.111$ [$-1.125$, $-1.097$] & [$-1.171$, $-1.033$] & [$-1.171$, $-1.033$] & [$-1.115$, $-1.115$] & never & $>$3 & -- \\
win95pts & semi-s. & Prt\allowbreak{}Mem $\to$ Problem5 & 9 & 1 & $\phantom{-}2.084$ & $\phantom{-}2.069$ [$\phantom{-}2.046$, $\phantom{-}2.093$] & [$\phantom{-}2.006$, $\phantom{-}2.169$] & [$\phantom{-}2.006$, $\phantom{-}2.169$] & [$\phantom{-}2.084$, $\phantom{-}2.084$] & never & $>$3 & -- \\
win95pts & semi-s. & Prt\allowbreak{}On $\to$ Prt\allowbreak{}Stat\allowbreak{}Off & 9 & 1 & $\phantom{-}1.011$ & $\phantom{-}1.012$ [$\phantom{-}0.998$, $\phantom{-}1.026$] & [$\phantom{-}0.000$, $\phantom{-}1.026$] & [$\phantom{-}0.000$, $\phantom{-}1.026$] & [$\phantom{-}0.000$, $\phantom{-}1.011$] & 1 & 1 & PrtOn$\to$PrtStatOff \\
win95pts & semi-s. & Prt\allowbreak{}PScript $\to$ Problem6 & 9 & 3 & $-1.881$ & $-1.892$ [$-1.922$, $-1.862$] & [$-2.183$, $-1.099$] & [$-2.183$, $-1.099$] & [$-2.143$, $-1.141$] & never & 1 & -- \\
win95pts & semi-s. & Prt\allowbreak{}Spool $\to$ GDIIN & 9 & 2 & $\phantom{-}1.037$ & $\phantom{-}1.042$ [$\phantom{-}1.028$, $\phantom{-}1.056$] & [$\phantom{-}0.848$, $\phantom{-}1.114$] & [$\phantom{-}0.848$, $\phantom{-}1.114$] & [$\phantom{-}1.037$, $\phantom{-}1.037$] & never & $>$3 & -- \\
win95pts & semi-s. & Prt\allowbreak{}Spool $\to$ Problem2 & 9 & 2 & $\phantom{-}1.219$ & $\phantom{-}1.208$ [$\phantom{-}1.184$, $\phantom{-}1.232$] & [$\phantom{-}0.232$, $\phantom{-}1.239$] & [$\phantom{-}0.232$, $\phantom{-}1.239$] & [$\phantom{-}0.328$, $\phantom{-}1.219$] & never & 1 & -- \\
win95pts & semi-s. & Tnr\allowbreak{}Spply $\to$ Prt\allowbreak{}Stat\allowbreak{}Toner & 9 & 1 & $\phantom{-}0.647$ & $\phantom{-}0.643$ [$\phantom{-}0.630$, $\phantom{-}0.657$] & [$\phantom{-}0.000$, $\phantom{-}0.657$] & [$\phantom{-}0.000$, $\phantom{-}0.657$] & [$\phantom{-}0.000$, $\phantom{-}0.647$] & 1 & 1 & TnrSpply$\to$PrtStatToner \\
\end{longtable}
\endgroup
\end{landscape}


\clearpage
\subsection{The 234 rows with no adjustment set}

Under the same condition, 234 of the 537 rows return no set at all: on the analyst's own graph the
outcome cannot descend from the treatment, so the effect is zero by construction and there is nothing
for a sensitivity report to be about. They are listed here by network so that the 537 rows of the
sweep are fully accounted for: $303 + 234 = 537$.

\begingroup\footnotesize
\setlength{\tabcolsep}{5pt}
\begin{longtable}{@{}lrl@{}}
\toprule
network & queries & verdict on the analyst's own graph \\
\midrule
\endfirsthead
\toprule
network & queries & verdict on the analyst's own graph \\
\midrule
\endhead
\bottomrule
\endlastfoot
Acid\_1996 & 1 & the outcome cannot descend from the treatment \\
Didelez\_2010 & 6 & the outcome cannot descend from the treatment \\
Kampen\_2014 & 10 & the outcome cannot descend from the treatment \\
Polzer\_2012 & 5 & the outcome cannot descend from the treatment \\
Schipf\_2010 & 8 & the outcome cannot descend from the treatment \\
Sebastiani\_2005 & 10 & the outcome cannot descend from the treatment \\
Shrier\_2008 & 7 & the outcome cannot descend from the treatment \\
Thoemmes\_2013 & 1 & the outcome cannot descend from the treatment \\
alarm & 5 & the outcome cannot descend from the treatment \\
andes & 10 & the outcome cannot descend from the treatment \\
arth150 & 18 & the outcome cannot descend from the treatment \\
asia & 7 & the outcome cannot descend from the treatment \\
barley & 6 & the outcome cannot descend from the treatment \\
child & 13 & the outcome cannot descend from the treatment \\
diabetes & 1 & the outcome cannot descend from the treatment \\
ecoli70 & 12 & the outcome cannot descend from the treatment \\
hailfinder & 17 & the outcome cannot descend from the treatment \\
hepar2 & 10 & the outcome cannot descend from the treatment \\
insurance & 11 & the outcome cannot descend from the treatment \\
magic-irri & 7 & the outcome cannot descend from the treatment \\
magic-niab & 6 & the outcome cannot descend from the treatment \\
mediator & 6 & the outcome cannot descend from the treatment \\
munin1 & 15 & the outcome cannot descend from the treatment \\
paths & 9 & the outcome cannot descend from the treatment \\
sachs & 13 & the outcome cannot descend from the treatment \\
water & 10 & the outcome cannot descend from the treatment \\
win95pts & 10 & the outcome cannot descend from the treatment \\
\midrule total & 234 & \\
\end{longtable}
\endgroup

